\begin{chapterpage}{Probabilitas}
  \chaptertitle{Probabilitas}
  \label{probability}
  \label{ch_probability}
  \chaptersection{basicsOfProbability}
  \chaptersection{conditionalProbabilitySection}
  \chaptersection{smallPop}
  \chaptersection{randomVariablesSection}
  \chaptersection{contDist}
\end{chapterpage}
\renewcommand{\chapterfolder}{ch_probability}

\index{probability|(}

\chapterintro{Probabilitas merupakan landasan statistika,
  dan Anda mungkin \mbox{sudah}
  mengenal banyak gagasan yang disajikan dalam bab ini.
  Namun, formalisasi konsep probabilitas kemungkinan besar masih
  baru bagi sebagian besar pembaca. \\

  \noindent%
  Walaupun bab ini memberikan landasan teoretis
  bagi gagasan-gagasan pada bab-bab selanjutnya dan membuka jalan
  menuju pemahaman yang lebih mendalam,
  penguasaan konsep-konsep yang diperkenalkan di bab ini
  tidak diperlukan untuk menerapkan
  metode yang diperkenalkan di bagian lain buku ini.}

%  This chapter provides a theoretical foundation for
%  the ideas introduced in later chapters.
%  However, this chapter is not strictly required to
%  understand or apply the methods introduced in the
%  rest of this book.}




\section{Mendefinisikan probabilitas}
\label{basicsOfProbability}

Statistika didasarkan pada probabilitas,
dan walaupun probabilitas tidak diperlukan untuk teknik-teknik
terapan dalam buku ini, pemahaman probabilitas dapat membantu Anda
memahami metode secara lebih mendalam dan membangun landasan yang lebih baik
untuk mata kuliah selanjutnya.


\subsection{Contoh pengantar}

Sebelum membahas gagasan teknis, mari kita lihat
beberapa contoh dasar yang mungkin terasa lebih akrab.

\begin{examplewrap}
\begin{nexample}{“Dadu”, bentuk tunggal dari dice, adalah kubus dengan enam sisi bernomor \resp{1}, \resp{2}, \resp{3}, \resp{4}, \resp{5}, dan \resp{6}. Berapa peluang memperoleh \resp{1} saat melempar dadu?}\label{probOf1}
Jika dadunya adil, peluang memperoleh \resp{1} sama dengan peluang memperoleh bilangan lain. Karena ada enam hasil, peluangnya adalah 1 banding 6 atau, secara ekuivalen, $1/6$.
\end{nexample}
\end{examplewrap}

\begin{examplewrap}
\begin{nexample}{Berapa peluang memperoleh \resp{1} atau \resp{2} pada lemparan berikutnya?}\label{probOf1Or2}
\resp{1} dan \resp{2} merupakan dua dari enam hasil yang mungkin dan sama-sama berpeluang, sehingga peluang memperoleh salah satu dari kedua hasil ini adalah $2/6 = 1/3$.
\end{nexample}
\end{examplewrap}

\begin{examplewrap}
\begin{nexample}{Berapa peluang memperoleh \resp{1}, \resp{2}, \resp{3}, \resp{4}, \resp{5}, atau \resp{6} pada lemparan berikutnya?}\label{probOf123456}
100\%. Hasilnya pasti salah satu dari bilangan tersebut.
\end{nexample}
\end{examplewrap}

\begin{examplewrap}
\begin{nexample}{Berapa peluang tidak memperoleh \resp{2}?}\label{probNot2}
Karena peluang memperoleh \resp{2} adalah $1/6$ atau $16.\bar{6}\%$, peluang tidak memperoleh \resp{2} adalah $100\% - 16.\bar{6}\%=83.\bar{3}\%$ atau $5/6$.

Sebagai alternatif, kita dapat menyadari bahwa tidak memperoleh \resp{2} sama dengan memperoleh \resp{1}, \resp{3}, \resp{4}, \resp{5}, atau \resp{6}. Kelima hasil itu mencakup lima dari enam hasil yang sama-sama berpeluang, sehingga probabilitasnya $5/6$.
\end{nexample}
\end{examplewrap}

\begin{examplewrap}
\begin{nexample}{Misalkan kita melempar dua dadu. Jika pada $1/6$ bagian dari seluruh lemparan dadu pertama menunjukkan \resp{1}, dan pada $1/6$ bagian dari waktu tersebut dadu kedua juga menunjukkan \resp{1}, berapa peluang memperoleh dua \resp{1}?}\label{probOf2Ones}
Jika dadu pertama menunjukkan \resp{1} pada $16.\bar{6}$\% lemparan dan pada $1/6$ bagian dari \emph{lemparan tersebut} dadu kedua juga menunjukkan \resp{1}, maka peluang kedua dadu menunjukkan \resp{1} adalah $(1/6)\times (1/6)$ atau $1/36$.
\end{nexample}
\end{examplewrap}


\D{\newpage}

\subsection{Probabilitas}

\index{random process|(}

Kita menggunakan probabilitas untuk membangun alat guna menggambarkan dan memahami keacakan yang tampak. Kita sering merumuskan probabilitas melalui \term{proses acak} yang menghasilkan suatu \term{hasil}.
\begin{center}
\begin{tabular}{lll}
Lempar dadu &$\rightarrow$ & \resp{1}, \resp{2}, \resp{3}, \resp{4}, \resp{5}, atau \resp{6} \\
Lempar koin &$\rightarrow$ & \resp{H} atau \resp{T} \\
\end{tabular}
\end{center}
Melempar dadu atau koin merupakan proses yang tampak acak, dan masing-masing menghasilkan suatu hasil.

\begin{onebox}{Probabilitas}
\term{Probabilitas} suatu hasil adalah proporsi kemunculan hasil tersebut jika kita mengamati proses acak itu sebanyak tak berhingga kali.
\end{onebox}

Probabilitas didefinisikan sebagai proporsi dan selalu bernilai antara 0~dan~1, termasuk kedua batasnya. Probabilitas juga dapat ditampilkan sebagai persentase antara 0\% dan 100\%.

Probabilitas dapat diilustrasikan dengan melempar dadu berkali-kali. Misalkan $\hat{p}_n$ adalah proporsi hasil yang berupa \resp{1} setelah $n$ lemparan pertama. Ketika jumlah lemparan bertambah, $\hat{p}_n$ akan konvergen ke probabilitas memperoleh \resp{1}, yaitu $p = 1/6$. Gambar~\ref{dieProp} menunjukkan konvergensi ini untuk 100.000 lemparan dadu. Kecenderungan $\hat{p}_n$ untuk stabil di sekitar $p$ dijelaskan oleh \term{Hukum Bilangan Besar}.

\begin{figure}[h]
\centering
\Figure[Diagram garis ditampilkan. Sumbu horizontal adalah "n (jumlah lemparan)", dengan nilai yang meningkat secara eksponensial dari 1 ke 10, 100, 1.000, 10.000, lalu 100.000. Sumbu vertikal adalah "p-hat sub n" dengan rentang 0,0 hingga sekitar 0,35. Garis putus-putus horizontal juga ditampilkan pada nilai satu per enam. Garis yang menunjukkan fraksi lemparan bernilai 1 mulai dari 0 pada lemparan pertama dan tetap di sana sampai sekitar lemparan ke-4, lalu melonjak ke 0,25 dan berosilasi, kemudian mencapai sekitar 0,35 pada 10 lemparan sebelum turun mendekati satu per enam. Sampai 100 lemparan, garis berosilasi antara 0,13 dan 0,22; dengan semakin banyak lemparan, garis makin stabil di sekitar satu per enam, menyimpang tidak lebih dari sekitar 0,03 sampai 1.000 lemparan. Sampai sekitar 5.000 lemparan penyimpangannya tidak lebih dari sekitar 0,015, dan setelah itu hampir tidak dapat dibedakan dari satu per enam selama lebih dari 5.000 lemparan.]{0.85}{dieProp}
\caption{Fraksi lemparan dadu yang menghasilkan 1 pada setiap tahap simulasi. Seiring bertambahnya jumlah lemparan, proporsinya cenderung mendekati probabilitas $1/6 \approx 0.167$.}
\label{dieProp}
\end{figure}

\begin{onebox}{Hukum Bilangan Besar}
Semakin banyak pengamatan dikumpulkan, proporsi $\hat{p}_n$ kemunculan suatu hasil tertentu akan konvergen ke probabilitas $p$ hasil tersebut.
\end{onebox}

Sesekali proporsi akan menyimpang dari probabilitas dan tampak menentang Hukum Bilangan Besar, seperti yang sering terjadi pada $\hat{p}_n$ di Gambar~\ref{dieProp}. Namun, penyimpangan itu mengecil ketika jumlah lemparan bertambah.

Di atas kita menulis $p$ sebagai probabilitas memperoleh \resp{1}. Probabilitas ini juga dapat ditulis sebagai
\begin{align*}
P(\text{memperoleh \resp{1}})
\end{align*}
Seiring terbiasanya kita dengan notasi ini, kita dapat menyingkatnya lebih lanjut. Misalnya, jika prosesnya jelas berupa “melempar dadu”, kita dapat menyingkat $P($memperoleh \resp{1}$)$ menjadi~$P($\resp{1}$)$.

\begin{exercisewrap}
\begin{nexercise} \label{randomProcessExercise}
Contoh proses acak adalah melempar dadu dan melempar koin. (a) Pikirkan proses acak lain. (b) Uraikan semua hasil yang mungkin dari proses itu. Sebagai contoh, melempar dadu adalah proses acak dengan hasil yang mungkin \mbox{\resp{1}, \resp{2}, ..., \resp{6}}.\footnotemark
\end{nexercise}
\end{exercisewrap}
\footnotetext{Berikut empat contoh. (i) Apakah seseorang akan sakit bulan depan atau tidak merupakan proses yang tampak acak, dengan hasil \resp{sakit} dan \resp{tidak}. (ii) Kita dapat \emph{membangkitkan} proses acak dengan memilih seseorang secara acak lalu mengukur tinggi badannya. Hasil proses ini berupa bilangan positif. (iii) Apakah pasar saham naik atau turun minggu depan merupakan proses yang tampak acak, dengan hasil yang mungkin \resp{naik}, \resp{turun}, dan \resp{tidak\us{}berubah}. Sebagai alternatif, perubahan persentase pasar saham dapat digunakan sebagai hasil numerik. (iv) Apakah teman sekamar Anda mencuci piring malam ini juga mungkin tampak seperti proses acak, dengan hasil \resp{mencuci\us{}piring} dan \resp{meninggalkan\us{}piring}.}

Hal-hal yang kita anggap sebagai proses acak belum tentu benar-benar acak; mungkin saja proses itu terlalu sulit dipahami secara tepat. Contoh keempat dalam jawaban catatan kaki Latihan Terbimbing~\ref{randomProcessExercise} mengisyaratkan bahwa perilaku teman sekamar adalah proses acak. Namun, meskipun perilaku teman sekamar tidak benar-benar acak, memodelkannya sebagai proses acak tetap dapat berguna.

%\begin{tipBox}{\tipBoxTitle{Modeling a process as random}
%It can be helpful to model a process as random even if it is not truly random.}
%\end{tipBox}

\index{random process|)}

\subsection{Hasil saling lepas atau saling eksklusif}

\index{disjoint|(}
\index{mutually exclusive|(}

Dua hasil disebut \term{saling lepas} atau \term{saling eksklusif} jika keduanya tidak mungkin terjadi bersamaan. Misalnya, jika kita melempar dadu, hasil \resp{1} dan \resp{2} saling lepas karena keduanya tidak mungkin muncul sekaligus. Sebaliknya, hasil \resp{1} dan “memperoleh bilangan ganjil” tidak saling lepas karena keduanya terjadi jika hasil lemparan adalah \resp{1}. Istilah \emph{saling lepas} dan \emph{saling eksklusif} ekuivalen dan dapat dipertukarkan.

Menghitung probabilitas hasil yang saling lepas cukup mudah. Saat melempar dadu, hasil \resp{1} dan \resp{2} saling lepas, sehingga probabilitas salah satu hasil itu terjadi dihitung dengan menjumlahkan probabilitas masing-masing:
\begin{align*}
P(\text{\resp{1} atau \resp{2}})
  = P(\text{\resp{1}})+P(\text{\resp{2}})
  = 1/6 + 1/6
  = 1/3
\end{align*}
Bagaimana dengan probabilitas memperoleh \resp{1}, \resp{2}, \resp{3}, \resp{4}, \resp{5}, atau \resp{6}? Sekali lagi, semua hasil tersebut saling lepas, sehingga probabilitasnya dijumlahkan:
\begin{align*}
&P(\text{\resp{1} atau \resp{2} atau \resp{3} atau \resp{4}
    atau \resp{5} atau \resp{6}}) \\
  &\quad = P(\text{\resp{1}})+P(\text{\resp{2}})
      + P(\text{\resp{3}})+P(\text{\resp{4}})
      + P(\text{\resp{5}})+P(\text{\resp{6}}) \\
  &\quad = 1/6 + 1/6 + 1/6 + 1/6 + 1/6 + 1/6
  = 1
\end{align*}
\term{Aturan Penjumlahan} menjamin ketepatan pendekatan ini ketika hasil-hasilnya saling lepas.

\begin{onebox}{Aturan Penjumlahan untuk hasil yang saling lepas}
  Jika $A_1$ dan $A_2$ menyatakan dua hasil yang saling lepas,
  probabilitas salah satunya terjadi diberikan oleh
  \begin{align*}
  P(A_1\text{ atau } A_2) = P(A_1) + P(A_2)
  \end{align*}
  Jika terdapat banyak hasil yang saling lepas, $A_1$, ..., $A_k$,
  probabilitas salah satu hasil tersebut terjadi adalah
  \begin{align*}
  P(A_1) + P(A_2) + \cdots + P(A_k)
  \end{align*}
\end{onebox}

\D{\newpage}

\begin{exercisewrap}
\begin{nexercise}
Kita tertarik pada probabilitas memperoleh \resp{1}, \resp{4}, atau \resp{5}. (a) Jelaskan mengapa hasil \resp{1}, \resp{4}, dan \resp{5} saling lepas. (b) Terapkan Aturan Penjumlahan untuk hasil yang saling lepas guna menentukan $P($\resp{1} atau \resp{4} atau \resp{5}$)$.\footnotemark
\end{nexercise}
\end{exercisewrap}
\footnotetext{(a) Proses acaknya adalah lemparan dadu, dan paling banyak satu dari hasil tersebut yang dapat muncul. Jadi, hasil-hasil itu saling lepas. (b)~$P($\resp{1} atau \resp{4} atau \resp{5}$) = P($\resp{1}$)+P($\resp{4}$)+P($\resp{5}$) = \frac{1}{6} + \frac{1}{6} + \frac{1}{6} = \frac{3}{6} = \frac{1}{2}$}

\index{data!loans|(}
\begin{exercisewrap}
\begin{nexercise}
Dalam himpunan data \data{loans} pada Bab~\ref{ch_summarizing_data},
variabel \var{homeownership} menjelaskan apakah peminjam
menyewa, memiliki hipotek, atau memiliki rumah sendiri.
Dari 10.000 peminjam, 3858 menyewa, 4789 memiliki hipotek,
dan 1353 memiliki rumah sendiri.\footnotemark{}
\begin{enumerate}[(a)]
\setlength{\itemsep}{0mm}
\item
    Apakah hasil \resp{rent}, \resp{mortgage},
    dan \resp{own} saling lepas?
\item
    Tentukan proporsi pinjaman dengan nilai \resp{mortgage}
    dan \resp{own} secara terpisah.
\item
    Gunakan Aturan Penjumlahan untuk hasil yang saling lepas
    guna menghitung probabilitas bahwa pinjaman yang dipilih
    secara acak dari himpunan data tersebut adalah milik seseorang
    yang memiliki hipotek atau rumah sendiri.
\end{enumerate}
\end{nexercise}
\end{exercisewrap}
\footnotetext{(a)~Ya. Setiap pinjaman dikategorikan ke tepat satu
level \var{homeownership}.
(b)~Hipotek: $\frac{4789}{10000} = 0.479$.
Milik sendiri: $\frac{1353}{10000} = 0.135$.
(c)~$P($\resp{mortgage} atau \resp{own}$) = P($\resp{mortgage}$) + P($\resp{own}$) = 0.479 + 0.135 = 0.614$.}
\index{data!loans|)}

\index{event|(}

Ilmuwan data jarang bekerja dengan hasil individual; mereka lebih sering mempertimbangkan \indexthis{\emph{himpunan}}{sets} atau \indexthis{\emph{kumpulan}}{collections} hasil. Misalkan $A$ menyatakan peristiwa ketika lemparan dadu menghasilkan \resp{1} atau \resp{2}, dan $B$ menyatakan peristiwa ketika lemparan menghasilkan \resp{4} atau \resp{6}. Kita menulis $A$ sebagai himpunan hasil $\{$\resp{1},~\resp{2}$\}$ dan $B=\{$\resp{4}, \resp{6}$\}$. Himpunan-himpunan ini lazim disebut \termsub{peristiwa}{event}. Karena $A$ dan $B$ tidak memiliki elemen yang sama, keduanya adalah peristiwa yang saling lepas. $A$ dan $B$ ditunjukkan pada Gambar~\ref{disjointSets}.

\begin{figure}[hhh]
  \centering
  \Figure[Enam bilangan ditampilkan berurutan: 1, 2, 3, 4, 5, dan 6. Bilangan 1 dan 2 dilingkari dan diberi label huruf “A”; bilangan 2 dan 3 dilingkari dan diberi label huruf “B”; bilangan 4 dan 6 dilingkari dengan label huruf “C”. (Lingkaran terakhir bukan lingkaran sebenarnya, melainkan kurungan yang tidak mencakup bilangan 5.)]{0.45}{disjointSets}
  \caption{Tiga peristiwa, $A$, $B$, dan $D$, terdiri atas
      hasil dari lemparan dadu.
      $A$ dan $B$ saling lepas karena tidak memiliki
      hasil yang sama.}
  \label{disjointSets}
\end{figure}

Aturan Penjumlahan berlaku untuk hasil maupun peristiwa yang saling lepas. Probabilitas salah satu peristiwa saling lepas $A$ atau $B$ terjadi adalah jumlah probabilitas masing-masing:
\begin{align*}
P(A\text{ atau }B) = P(A) + P(B) = 1/3 + 1/3 = 2/3
\end{align*}

\begin{exercisewrap}
\begin{nexercise}
(a) Verifikasi bahwa probabilitas peristiwa $A$, $P(A)$,
adalah $1/3$ menggunakan Aturan Penjumlahan.
(b)~Lakukan hal yang sama untuk peristiwa $B$.\footnotemark
\end{nexercise}
\end{exercisewrap}
\footnotetext{(a)~$P(A) = P($\resp{1} atau \resp{2}$)
    = P($\resp{1}$) + P($\resp{2}$)
    = \frac{1}{6} + \frac{1}{6}
    = \frac{2}{6}
    = \frac{1}{3}$.
    (b)~Dengan cara serupa, $P(B) = 1/3$.}

\begin{exercisewrap}
\begin{nexercise} \label{exerExaminingDisjointSetsABD}
(a) Dengan menggunakan Gambar~\ref{disjointSets} sebagai acuan, hasil apa yang direpresentasikan oleh peristiwa $D$? (b) Apakah peristiwa $B$ dan $D$ saling lepas? (c) Apakah peristiwa $A$ dan $D$ saling lepas?\footnotemark
\end{nexercise}
\end{exercisewrap}
\footnotetext{(a)~Hasil \resp{2} dan \resp{3}. (b)~Ya, peristiwa $B$ dan $D$ saling lepas karena tidak memiliki hasil yang sama. (c)~Peristiwa $A$ dan $D$ memiliki hasil yang sama, yaitu \resp{2}, sehingga tidak saling lepas.}

\begin{exercisewrap}
\begin{nexercise}
Dalam Latihan Terbimbing~\ref{exerExaminingDisjointSetsABD}, Anda mengonfirmasi bahwa $B$ dan $D$ pada Gambar~\ref{disjointSets} saling lepas. Hitung probabilitas peristiwa $B$ atau $D$ terjadi.\footnotemark
\end{nexercise}
\end{exercisewrap}
\footnotetext{Karena $B$ dan $D$ adalah peristiwa yang saling lepas, gunakan Aturan Penjumlahan: $P(B$ atau $D) = P(B) + P(D) = \frac{1}{3} + \frac{1}{3} = \frac{2}{3}$.}

\index{event|)}
\index{disjoint|)}
\index{mutually exclusive|)}


\D{\newpage}

\subsection{Probabilitas ketika peristiwa tidak saling lepas}

Mari kita pertimbangkan perhitungan untuk dua peristiwa yang tidak saling lepas dalam konteks \indexthis{dek standar 52 kartu}{deck of cards}, yang ditampilkan pada Gambar~\ref{deckOfCards}. Jika Anda belum mengenal kartu dalam dek standar, lihat catatan kaki.\footnote{Ke-52 kartu dibagi menjadi empat \term{jenis}: $\clubsuit$ (keriting), {\color{redcards}$\diamondsuit$} (wajik), {\color{redcards}$\heartsuit$} (hati), dan $\spadesuit$ (sekop). Setiap jenis memiliki 13 kartu berlabel: \resp{2}, \resp{3}, ..., \resp{10}, \resp{J} (jack), \resp{Q} (queen), \resp{K} (king), dan \resp{A} (ace). Jadi, setiap kartu merupakan kombinasi unik antara jenis dan label, misalnya {\color{redcards}\resp{4$\heartsuit$}} dan \resp{J$\clubsuit$}. Ke-12 kartu berupa jack, queen, dan king disebut \termsub{\resp{kartu bergambar}}{face card}. Kartu berjenis {\color{redcards}$\diamondsuit$} atau {\color{redcards}$\heartsuit$} biasanya berwarna {\color{redcards}merah}, sedangkan dua jenis lainnya biasanya berwarna hitam.}

\begin{figure}[h]
\centering
\begin{tabular}{lll lll lll lll l}
\resp{2$\clubsuit$} & \resp{3$\clubsuit$} & \resp{4$\clubsuit$} & \resp{5$\clubsuit$} & \resp{6$\clubsuit$} & \resp{7$\clubsuit$} & \resp{8$\clubsuit$} & \resp{9$\clubsuit$} & \resp{10$\clubsuit$} & \resp{J$\clubsuit$} & \resp{Q$\clubsuit$} & \resp{K$\clubsuit$} & \resp{A$\clubsuit$}  \\
\color{redcards} \resp{2$\diamondsuit$} & \color{redcards}\resp{3$\diamondsuit$} & \color{redcards}\resp{4$\diamondsuit$} & \color{redcards}\resp{5$\diamondsuit$} & \color{redcards}\resp{6$\diamondsuit$} & \color{redcards}\resp{7$\diamondsuit$} & \color{redcards}\resp{8$\diamondsuit$} & \color{redcards}\resp{9$\diamondsuit$} & \color{redcards}\resp{10$\diamondsuit$} & \color{redcards}\resp{J$\diamondsuit$} & \color{redcards}\resp{Q$\diamondsuit$} & \color{redcards}\resp{K$\diamondsuit$} & \color{redcards}\resp{A$\diamondsuit$} \\
\color{redcards}\resp{2$\heartsuit$} & \color{redcards}\resp{3$\heartsuit$} & \color{redcards}\resp{4$\heartsuit$} & \color{redcards}\resp{5$\heartsuit$} & \color{redcards}\resp{6$\heartsuit$} & \color{redcards}\resp{7$\heartsuit$} & \color{redcards}\resp{8$\heartsuit$} & \color{redcards}\resp{9$\heartsuit$} & \color{redcards}\resp{10$\heartsuit$} & \color{redcards}\resp{J$\heartsuit$} & \color{redcards}\resp{Q$\heartsuit$} & \color{redcards}\resp{K$\heartsuit$} & \color{redcards}\resp{A$\heartsuit$} \\
\resp{2$\spadesuit$} & \resp{3$\spadesuit$} & \resp{4$\spadesuit$} & \resp{5$\spadesuit$} & \resp{6$\spadesuit$} & \resp{7$\spadesuit$} & \resp{8$\spadesuit$} & \resp{9$\spadesuit$} & \resp{10$\spadesuit$} & \resp{J$\spadesuit$} & \resp{Q$\spadesuit$} & \resp{K$\spadesuit$} & \resp{A$\spadesuit$}
\end{tabular}
\caption{Representasi 52 kartu unik dalam sebuah dek.}
\label{deckOfCards}
\end{figure}

\begin{exercisewrap}
\begin{nexercise}
(a) Berapa probabilitas kartu yang dipilih secara acak merupakan wajik? (b)~Berapa probabilitas kartu yang dipilih secara acak merupakan kartu bergambar?\footnotemark
\end{nexercise}
\end{exercisewrap}
\footnotetext{(a) Ada 52 kartu dan 13 kartu wajik. Jika kartu dikocok dengan baik, setiap kartu memiliki peluang yang sama untuk terambil, sehingga probabilitas kartu yang dipilih secara acak merupakan wajik adalah $P({\color{redcards}\diamondsuit}) = \frac{13}{52} = 0.250$. (b)~Demikian pula, ada 12 kartu bergambar, sehingga $P($kartu bergambar$) = \frac{12}{52} = \frac{3}{13} = 0.231$.}

\term{Diagram Venn} berguna ketika hasil dapat dikategorikan sebagai “masuk” atau “tidak masuk” untuk dua atau tiga variabel, atribut, atau proses acak. Diagram Venn pada Gambar~\ref{cardsDiamondFaceVenn} menggunakan satu lingkaran untuk mewakili wajik dan lingkaran lain untuk mewakili kartu bergambar. Jika sebuah kartu merupakan wajik sekaligus kartu bergambar, kartu itu berada di irisan kedua lingkaran. Jika kartu itu wajik tetapi bukan kartu bergambar, kartu itu berada di bagian lingkaran kiri yang tidak termasuk lingkaran kanan (dan seterusnya). Jumlah kartu wajik diberikan oleh seluruh kartu di dalam lingkaran wajik: $10+3=13$. Probabilitasnya juga ditampilkan (misalnya, $10/52 = 0.1923$).

\begin{figure}[h]
\centering
\Figure[Diagram Venn ditampilkan. Satu lingkaran berlabel “Wajik” dengan proporsi total 0,25 dan lingkaran kedua berlabel “Kartu bergambar” dengan proporsi total 0,2308. Kedua lingkaran beririsan dan berbagi 3 kartu, dengan proporsi 0,0577 dari satu dek. Bagian lingkaran wajik yang tidak beririsan berlabel “10” untuk 10 kartu dan proporsi 0,1923. Bagian lingkaran kartu bergambar yang tidak beririsan berlabel “9” untuk 9 kartu dan proporsi 0,2308. Gambar juga mencatat bahwa terdapat 30 kartu yang bukan wajik maupun kartu bergambar.]{0.65}{cardsDiamondFaceVenn}
\caption{Diagram Venn untuk wajik dan kartu bergambar.}
\label{cardsDiamondFaceVenn}
\end{figure}

Misalkan $A$ menyatakan peristiwa bahwa kartu yang dipilih secara acak adalah wajik dan $B$ menyatakan peristiwa bahwa kartu tersebut merupakan kartu bergambar. Bagaimana menghitung $P(A$ atau $B)$? Peristiwa $A$ dan $B$ tidak saling lepas--kartu {\color{redcards}$J\diamondsuit$}, {\color{redcards}$Q\diamondsuit$}, dan {\color{redcards}$K\diamondsuit$} termasuk dalam kedua kategori--sehingga kita tidak dapat menggunakan Aturan Penjumlahan untuk peristiwa yang saling lepas. Sebagai gantinya, kita menggunakan diagram Venn. Kita mulai dengan menjumlahkan probabilitas kedua peristiwa:
\begin{align*}
P(A) + P(B)
  = P({\color{redcards}\diamondsuit}) + P(\text{kartu bergambar})
  = 13/52 + 12/52
\end{align*}

\D{\newpage}

\noindent%
Namun, tiga kartu yang termasuk dalam kedua peristiwa telah dihitung dua kali, sekali pada setiap probabilitas. Kita harus memperbaiki penghitungan ganda ini:
\begin{align*}
P(A\text{ atau } B)
  &= P({\color{redcards}\diamondsuit}\text{ atau kartu bergambar}) \\
  &= P({\color{redcards}\diamondsuit}) + P(\text{kartu bergambar})
      - P({\color{redcards}\diamondsuit}\text{ dan kartu bergambar}) \\
  &= 13/52 + 12/52 - 3/52 \\
  &= 22/52 = 11/26
\end{align*}
Persamaan ini merupakan contoh \term{Aturan Penjumlahan Umum}.

\begin{onebox}{Aturan Penjumlahan Umum}
  Jika $A$ dan $B$ adalah sembarang dua peristiwa, baik saling lepas maupun tidak,
  probabilitas bahwa setidaknya salah satunya terjadi adalah
  \begin{align*}
  P(A\text{ atau }B) = P(A) + P(B) - P(A\text{ dan }B)
  \end{align*}
  dengan $P(A$ dan $B)$ adalah probabilitas kedua peristiwa terjadi.
\end{onebox}

\begin{tipBox}{\tipBoxTitle{“atau” bersifat inklusif}
Ketika kita menulis “atau” dalam statistika, yang dimaksud adalah “dan/atau” kecuali dinyatakan secara eksplisit. Jadi, $A$ atau $B$ terjadi berarti $A$, $B$, atau keduanya, $A$ dan $B$, terjadi.}
\end{tipBox}

\begin{exercisewrap}
\begin{nexercise}
(a) Jika $A$ dan $B$ saling lepas, jelaskan mengapa hal ini berarti $P(A$ dan $B) = 0$. (b) Dengan menggunakan bagian (a), verifikasi bahwa Aturan Penjumlahan Umum menyederhana menjadi Aturan Penjumlahan untuk peristiwa saling lepas jika $A$ dan $B$ saling lepas.\footnotemark
\end{nexercise}
\end{exercisewrap}
\footnotetext{(a) Jika $A$ dan $B$ saling lepas, $A$ dan $B$ tidak mungkin terjadi secara bersamaan. (b) Jika $A$ dan $B$ saling lepas, suku terakhir $P(A\text{ dan }B)$ dalam rumus Aturan Penjumlahan Umum bernilai 0 (lihat bagian (a)), sehingga tersisa Aturan Penjumlahan untuk peristiwa saling lepas.}

\index{data!loans|(}

\begin{exercisewrap}
\begin{nexercise}\label{emailSpamNumberVennExer}
% library(openintro); d <- loans_full_schema; table(d[,c("application_type", "homeownership")]); table(d[,c("application_type")]); table(d[,c("homeownership")])
Dalam himpunan data \data{loans} yang memuat 10.000 pinjaman,
1495 pinjaman berasal dari permohonan bersama
(misalnya, pasangan mengajukan bersama),
4789 pemohon memiliki hipotek,
dan 950 memiliki kedua karakteristik tersebut.
Buat diagram Venn untuk situasi ini.\footnotemark
\end{nexercise}
\end{exercisewrap}
\footnotetext{%
  \begin{minipage}[t]{0.65\textwidth}
  Baik jumlah maupun {\color{oiB}probabilitas} yang bersesuaian
  (misalnya, $3839/10000 = 0.384$) ditampilkan.
  Perhatikan bahwa jumlah pinjaman pada lingkaran kiri
  adalah $3839 + 950 = 4789$, sedangkan jumlah
  pada lingkaran kanan adalah $950 + 545 = 1495$.
  \end{minipage}\ %
  \begin{minipage}[c]{0.3\textwidth}
  \hfill\Figure[Diagram Venn dengan dua lingkaran ditampilkan. Lingkaran pertama berlabel “pemohon memiliki hipotek” dan lingkaran kedua berlabel “permohonan bersama”; kedua lingkaran sebagian beririsan. Pada lingkaran pertama, bagian yang tidak beririsan menunjukkan jumlah 3839 dan proporsi 0,384. Bagian lingkaran permohonan bersama yang tidak beririsan berlabel 545 dengan proporsi 0,055. Bagian irisan berlabel jumlah 950 dan proporsi 0,095. Di luar kedua lingkaran dicatat bahwa “pinjaman lain” berjumlah 10.000 dikurangi 3.839, 950, dan 545, yaitu 4666 dengan proporsi 0,467.]{}{loans_app_type_home_venn} \vspace{-13mm}
  \end{minipage}}

\begin{exercisewrap}
\begin{nexercise}
(a)~Gunakan diagram Venn dari Latihan Terbimbing~\ref{emailSpamNumberVennExer} untuk menentukan
probabilitas bahwa pinjaman yang dipilih secara acak dari himpunan data \data{loans}
berasal dari permohonan bersama dengan pasangan yang memiliki
hipotek.
(b)~Berapa probabilitas pinjaman memiliki salah satu
atribut tersebut?\footnotemark
\end{nexercise}
\end{exercisewrap}
\footnotetext{%
  (a)~Jawabannya direpresentasikan oleh irisan
  kedua lingkaran: 0.095.
  (b)~Ini adalah jumlah tiga probabilitas yang saling lepas yang ditunjukkan
  di dalam lingkaran: $0.384 + 0.095 + 0.055 = 0.534$
  (berbeda 0.001 karena pembulatan).}

\index{data!loans|)}


\D{\newpage}

\subsection{Distribusi probabilitas}

\termsub{Distribusi probabilitas}{probability!distribution} adalah tabel semua hasil yang saling lepas beserta probabilitasnya. Gambar~\ref{diceProb} menunjukkan distribusi probabilitas jumlah dua dadu.

\begin{figure}[h] \small
\centering
\begin{tabular}{l ccc ccc ccc cc}
  \hline
  \ \vspace{-3mm} \\
Jumlah dadu\vspace{0.3mm} & 2 & 3 & 4 & 5 & 6 & 7 & 8 & 9 & 10 & 11 & 12  \\
Probabilitas & $\frac{1}{36}$ & $\frac{2}{36}$ & $\frac{3}{36}$ & $\frac{4}{36}$ & $\frac{5}{36}$ & $\frac{6}{36}$ & $\frac{5}{36}$ & $\frac{4}{36}$ & $\frac{3}{36}$ & $\frac{2}{36}$ & $\frac{1}{36}$\vspace{1mm} \\
   \hline
\end{tabular}
\caption{Distribusi probabilitas jumlah dua dadu.}
\label{diceProb}
\end{figure}

\begin{onebox}{Aturan distribusi probabilitas}
Distribusi probabilitas adalah daftar hasil yang mungkin beserta probabilitasnya, yang memenuhi tiga aturan berikut: \vspace{-2mm}
\begin{enumerate}
\setlength{\itemsep}{0mm}
\item Hasil yang tercantum harus saling lepas.
\item Setiap probabilitas harus berada antara 0 dan 1.
\item Jumlah seluruh probabilitas harus 1. \vspace{1mm}
\end{enumerate}
\end{onebox}

\begin{exercisewrap}
\begin{nexercise}\label{usHouseholdIncomeDistsExercise}
Gambar~\ref{usHouseholdIncomeDists} menyajikan tiga distribusi pendapatan rumah tangga di Amerika Serikat. Hanya satu yang benar. Yang mana? Apa kesalahan pada dua lainnya?\footnotemark
\end{nexercise}
\end{exercisewrap}
\footnotetext{Probabilitas pada (a) tidak berjumlah~1.
    Probabilitas kedua pada (b) bernilai negatif.
    Dengan demikian tersisa (c), yang memang memenuhi
    syarat sebuah distribusi.
    Salah satu dari ketiganya disebut sebagai
    distribusi pendapatan rumah tangga AS yang sebenarnya,
    sehingga jawabannya pasti (c).}

\begin{figure}[h]
\centering
\begin{tabular}{r | cc cc}
  \hline
Rentang pendapatan & \$0-25k & \$25k-50k & \$50k-100k & \$100k+ \\
  \hline
(a)\hspace{0.2mm}	 & 0.18 & 0.39 & 0.33 & 0.16 \\
(b)				 & 0.38 & -0.27 & 0.52 & 0.37 \\
(c)\hspace{0.2mm}	 & 0.28 & 0.27 & 0.29 & 0.16 \\
  \hline
\end{tabular}
\caption{Distribusi pendapatan rumah tangga AS yang diusulkan (Latihan Terbimbing~\ref{usHouseholdIncomeDistsExercise}).}
\label{usHouseholdIncomeDists}
\end{figure}

Bab~\ref{introductionToData} menekankan pentingnya membuat plot data untuk memperoleh ringkasan cepat. Distribusi probabilitas juga dapat diringkas dengan diagram batang. Misalnya, distribusi pendapatan rumah tangga AS ditampilkan sebagai diagram batang pada Gambar~\ref{usHouseholdIncomeDistBar}. %\footnote{Plot distribusi juga dapat dibuat ketika pendapatan tidak dibagi secara artifisial ke dalam empat kelompok. Distribusi \emph{kontinu} dibahas pada Bagian~\ref{contDist}.}
Distribusi probabilitas jumlah dua dadu ditampilkan pada Gambar~\ref{diceProb} dan diplot pada Gambar~\ref{diceSumDist}.

\begin{figure}[h]
  \centering
  \Figure[Diagram batang “Pendapatan Rumah Tangga AS” ditampilkan dengan empat kelompok pendapatan. Sumbu vertikal berlabel “Probabilitas”. Kelompok pertama adalah \$0 sampai \$25.000 dengan tinggi batang sekitar 0,28; kelompok kedua \$25.000 sampai \$50.000 dengan tinggi sekitar 0,27; kelompok ketiga \$50.000 sampai \$100.000 dengan tinggi sekitar 0,28; dan kelompok keempat di atas \$100.000 dengan tinggi sekitar 0,15.]{0.65}{usHouseholdIncomeDistBar}
  \caption{Distribusi probabilitas pendapatan rumah tangga AS.}
  \label{usHouseholdIncomeDistBar}
\end{figure}

\begin{figure}
  \centering
  \Figure[Diagram batang jumlah dua dadu ditampilkan; nilainya dapat berupa 2, 3, 4, 5, dan seterusnya sampai 12. Sumbu vertikal berlabel “Probabilitas”. Tinggi batang untuk 2 sekitar 0,025; 3 sebesar 0,055; 4 sebesar 0,09; 5 sebesar 0,115; 6 sebesar 0,14; 7 sebesar 0,165; 8 sebesar 0,14; 9 sebesar 0,115; 10 sebesar 0,09; 11 sebesar 0,055; dan 12 sebesar 0,025.]{0.67}{diceSumDist}
  \caption{Distribusi probabilitas jumlah dua dadu.}
  \label{diceSumDist}
\end{figure}

Pada diagram batang ini, tinggi batang merepresentasikan probabilitas hasil. Jika hasilnya numerik dan diskret, biasanya (secara visual) nyaman membuat diagram batang yang menyerupai histogram, seperti pada jumlah dua dadu. Contoh lain penempatan batang pada lokasi masing-masing ditunjukkan pada Gambar~\ref{bookCostDist} di halaman~\pageref{bookCostDist}.

\subsection{Komplemen suatu peristiwa}

Melempar dadu menghasilkan nilai dalam himpunan $\{$\resp{1}, \resp{2}, \resp{3}, \resp{4}, \resp{5}, \resp{6}$\}$. Himpunan semua hasil yang mungkin ini disebut \term{ruang sampel} ($S$)\index{S@$S$} untuk lemparan dadu. Kita sering menggunakan ruang sampel untuk memeriksa situasi ketika suatu peristiwa tidak terjadi.

Misalkan $D=\{$\resp{2}, \resp{3}$\}$ menyatakan peristiwa bahwa hasil lemparan dadu adalah \resp{2} atau \resp{3}. Kemudian \term{komplemen} $D$ menyatakan semua hasil dalam ruang sampel yang bukan anggota $D$, dilambangkan dengan $D^c = \{$\resp{1}, \resp{4}, \resp{5}, \resp{6}$\}$. Dengan kata lain, $D^c$ adalah himpunan semua hasil yang mungkin yang belum termasuk dalam $D$. Gambar~\ref{complementOfD} menunjukkan hubungan antara $D$, $D^c$, dan ruang sampel $S$.

\begin{figure}[hht]
  \centering
  \Figure[Bilangan 1, 2, 3, 4, 5, dan 6 ditampilkan berurutan. Bilangan 2 dan 3 dilingkari serta diberi label “D”. Bilangan 1, 4, 5, dan 6 dilingkari serta diberi label “D-pangkat-C” sebagai komplemen D. Semua bilangan kemudian dilingkari lagi dengan kurungan lebih besar berlabel “S” sebagai ruang sampel.]{0.55}{complementOfD}
  \caption{Peristiwa $D=\{$\resp{2}, \resp{3}$\}$ dan komplemennya,
      $D^c = \{$\resp{1}, \resp{4}, \resp{5}, \resp{6}$\}$.
      $S$~menyatakan ruang sampel, yaitu himpunan
      semua hasil yang mungkin.}
  \label{complementOfD}
\end{figure}

\begin{exercisewrap}
\begin{nexercise}
(a) Hitung $P(D^c) = P($memperoleh \resp{1}, \resp{4}, \resp{5}, atau \resp{6}$)$. (b) Berapakah $P(D) + P(D^c)$?\footnotemark
\end{nexercise}
\end{exercisewrap}
\footnotetext{(a)~Hasil-hasilnya saling lepas dan masing-masing memiliki probabilitas $1/6$, sehingga probabilitas totalnya $4/6=2/3$. (b)~Kita juga dapat melihat bahwa $P(D)=\frac{1}{6} + \frac{1}{6} = 1/3$. Karena $D$ dan $D^c$ saling lepas, $P(D) + P(D^c) = 1$.}

\begin{exercisewrap}
\begin{nexercise}
Peristiwa $A=\{$\resp{1}, \resp{2}$\}$ dan $B=\{$\resp{4}, \resp{6}$\}$ ditunjukkan pada Gambar~\ref{disjointSets} di halaman~\pageref{disjointSets}. (a) Tuliskan apa yang direpresentasikan oleh $A^c$ dan $B^c$. (b)~Hitung $P(A^c)$ dan $P(B^c)$. (c)~Hitung $P(A)+P(A^c)$ dan $P(B)+P(B^c)$.\footnotemark
\end{nexercise}
\end{exercisewrap}
\footnotetext{Jawaban singkat: (a)~$A^c=\{$\resp{3}, \resp{4}, \resp{5}, \resp{6}$\}$ dan $B^c=\{$\resp{1}, \resp{2}, \resp{3}, \resp{5}$\}$. (b)~Karena setiap hasil saling lepas, jumlahkan probabilitas hasil individual untuk memperoleh $P(A^c)=2/3$ dan $P(B^c)=2/3$. (c)~$A$~dan~$A^c$ saling lepas, demikian pula $B$~dan~$B^c$. Oleh karena itu, $P(A) + P(A^c) = 1$ dan $P(B) + P(B^c) = 1$.}

\D{\newpage}

Komplemen suatu peristiwa $A$ dibentuk agar memiliki dua sifat penting: (i) setiap hasil yang mungkin dan bukan anggota $A$ merupakan anggota $A^c$, dan (ii) $A$ dan $A^c$ saling lepas. Sifat (i) menyiratkan
\begin{align*}
P(A\text{ atau }A^c) = 1
\end{align*}
Artinya, jika hasil bukan anggota $A$, hasil itu harus direpresentasikan dalam $A^c$. Kita menggunakan Aturan Penjumlahan untuk peristiwa saling lepas guna menerapkan Sifat (ii):
\begin{align*}
P(A\text{ atau }A^c) = P(A) + P(A^c)
\end{align*}
Menggabungkan dua persamaan terakhir menghasilkan hubungan yang sangat berguna antara probabilitas suatu peristiwa dan komplemennya.

\begin{onebox}{Komplemen}
  Komplemen peristiwa $A$ dilambangkan $A^c$, dan $A^c$
  menyatakan semua hasil yang bukan anggota~$A$. $A$ dan $A^c$
  berhubungan secara matematis:
  \begin{align*}
  P(A) + P(A^c) = 1, \quad\text{i.e.}\quad P(A) = 1-P(A^c)
  \end{align*}
\end{onebox}

Dalam contoh sederhana, menghitung $A$ atau $A^c$ dapat dilakukan dalam beberapa langkah. Namun, menggunakan komplemen dapat menghemat banyak waktu ketika masalah menjadi lebih kompleks.

\begin{exercisewrap}
\begin{nexercise}
Misalkan $A$ menyatakan peristiwa ketika kita melempar dua dadu dan jumlahnya kurang dari \resp{12}. (a)~Apa yang direpresentasikan oleh peristiwa $A^c$? (b)~Tentukan $P(A^c)$ dari Gambar~\ref{diceProb} pada halaman~\pageref{diceProb}. (c)~Tentukan $P(A)$.\footnotemark
\end{nexercise}
\end{exercisewrap}
\footnotetext{(a)~Komplemen $A$: jumlahnya sama dengan \resp{12}. (b)~$P(A^c) = 1/36$. (c)~Gunakan probabilitas komplemen dari bagian (b), $P(A^c) = 1/36$, dan persamaan komplemen: $P($kurang dari \resp{12}$) = 1 - P($\resp{12}$) = 1 - 1/36 = 35/36$.}

\begin{exercisewrap}
\begin{nexercise}
Tentukan probabilitas berikut untuk lemparan dua dadu:\footnotemark
\begin{enumerate}[(a)]
\setlength{\itemsep}{0mm}
\item Jumlah dadu \emph{bukan} \resp{6}.
\item Jumlahnya paling sedikit \resp{4}.
    Artinya, tentukan probabilitas peristiwa
    $B = \{$\resp{4}, \resp{5}, ..., \resp{12}$\}$.
\item Jumlahnya paling banyak \resp{10}.
    Artinya, tentukan probabilitas peristiwa
    $D=\{$\resp{2}, \resp{3}, ..., \resp{10}$\}$.
\end{enumerate}
\end{nexercise}
\end{exercisewrap}
\footnotetext{(a)~Pertama, cari $P($\resp{6}$)=5/36$, lalu gunakan komplemen: $P($bukan \resp{6}$) = 1 - P($\resp{6}$) = 31/36$.

(b)~Pertama, cari komplemennya, yang membutuhkan jauh lebih sedikit usaha: $P($\resp{2} atau \resp{3}$)=1/36+2/36=1/12$. Kemudian hitung $P(B) = 1-P(B^c) = 1-1/12 = 11/12$.

(c)~Seperti sebelumnya, mencari komplemen adalah cara cerdas untuk menentukan $P(D)$. Pertama, cari $P(D^c) = P($\resp{11} atau \resp{12}$)=2/36 + 1/36=1/12$. Kemudian hitung $P(D) = 1 - P(D^c) = 11/12$.}


\subsection{Kemandirian}
\label{probabilityIndependence}

Sebagaimana variabel dan pengamatan dapat saling independen, proses acak juga dapat saling independen. Dua proses \term{independen} jika mengetahui hasil salah satunya tidak memberikan informasi berguna tentang hasil yang lain. Misalnya, melempar koin dan melempar dadu adalah dua proses independen--mengetahui koin menunjukkan kepala tidak membantu menentukan hasil lemparan dadu. Sebaliknya, harga saham biasanya bergerak naik atau turun bersama-sama, sehingga tidak independen.

Contoh~\ref{probOf2Ones} memberikan contoh dasar dua proses independen: melempar dua dadu. Kita ingin menentukan probabilitas keduanya menunjukkan \resp{1}. Misalkan salah satu dadu berwarna merah dan yang lain putih. Jika hasil dadu merah adalah \resp{1}, hal itu tidak memberikan informasi tentang hasil dadu putih. Kita telah menjumpai pertanyaan yang sama pada Contoh~\ref{probOf2Ones} (halaman~\pageref{probOf2Ones}), ketika probabilitas dihitung dengan penalaran berikut: dadu merah menunjukkan \resp{1} pada $1/6$ bagian waktu, dan pada $1/6$ bagian dari \emph{waktu tersebut} dadu putih juga menunjukkan \resp{1}. Hal ini diilustrasikan pada Gambar~\ref{indepForRollingTwo1s}. Karena lemparan-lemparan itu independen, probabilitas hasil yang bersesuaian dapat dikalikan untuk memperoleh jawaban akhir: $(1/6)\times(1/6)=1/36$. Penalaran ini dapat digeneralisasi ke banyak proses independen.

\begin{figure}[hht]
\centering
\Figure[Persegi panjang hitam menguraikan gambar dan berlabel “Semua lemparan”. Di dalamnya, pita vertikal selebar kira-kira seperenam diarsir dan diberi label “seperenam lemparan pertama bernilai 1”. Bagian horizontal yang mewakili kira-kira seperenam dari pita vertikal itu diarsir berbeda dan diberi label “pada seperenam waktu ketika lemparan pertama bernilai 1, lemparan kedua juga bernilai 1”.]{0.6}{indepForRollingTwo1s}
\caption{Pada $1/6$ bagian waktu, lemparan pertama bernilai \resp{1}. Kemudian pada $1/6$ bagian dari \emph{waktu tersebut}, lemparan kedua juga bernilai \resp{1}.}
\label{indepForRollingTwo1s}
\end{figure}

\begin{examplewrap}
\begin{nexample}{Bagaimana jika ada dadu biru yang juga independen dari dua dadu lainnya? Berapa probabilitas melempar ketiga dadu dan semuanya menunjukkan \resp{1}?}\label{threeDice}
Logika yang sama seperti pada Contoh~\ref{probOf2Ones} berlaku. Jika dadu putih dan merah sama-sama \resp{1} pada $1/36$ bagian waktu, maka pada $1/6$ bagian dari \emph{waktu tersebut} dadu biru juga akan \resp{1}, sehingga kita mengalikan:
{\begin{align*}
P(putih=\text{\small\resp{1} dan } merah=\text{\small\resp{1} dan } biru=\text{\small\resp{1}})
	&= P(putih=\text{\small\resp{1}})\times P(merah=\text{\small\resp{1}})\times P(biru=\text{\small\resp{1}}) \\
	&= (1/6)\times (1/6)\times (1/6)
	= 1/216
\end{align*}} \vspace{-7mm}
\end{nexample}
\end{examplewrap}

Contoh~\ref{threeDice} mengilustrasikan apa yang disebut Aturan Perkalian untuk proses independen.

\begin{onebox}{Aturan Perkalian untuk proses independen}
  \index{Multiplication Rule|textbf}%
  Jika $A$ dan $B$ menyatakan peristiwa dari dua proses
  yang berbeda dan independen, probabilitas bahwa $A$
  dan $B$ terjadi dapat dihitung sebagai hasil kali
  probabilitas masing-masing:
  \begin{align*}
  P(A \text{ dan }B) = P(A) \times  P(B)
  \end{align*}
  Demikian pula, jika terdapat $k$ peristiwa $A_1$, ..., $A_k$
  dari $k$ proses independen, probabilitas semuanya terjadi adalah
  \begin{align*}
  P(A_1) \times  P(A_2)\times  \cdots \times  P(A_k)
  \end{align*}\vspace{-6mm}
\end{onebox}

\begin{exercisewrap}
\begin{nexercise} \label{ex2Handedness}
Sekitar 9\% orang kidal. Misalkan 2 orang dipilih secara acak dari populasi AS. Karena ukuran sampel 2 sangat kecil dibandingkan populasi, wajar untuk menganggap kedua orang ini independen. (a)~Berapa probabilitas keduanya kidal? (b)~Berapa probabilitas keduanya tidak kidal?\footnotemark
\end{nexercise}
\end{exercisewrap}
\footnotetext{(a) Probabilitas orang pertama kidal adalah $0.09$, sama dengan orang kedua. Kita menerapkan Aturan Perkalian untuk proses independen untuk menentukan probabilitas keduanya kidal: $0.09\times 0.09 = 0.0081$.

(b) Wajar untuk menganggap proporsi orang ambidekster (dapat menggunakan tangan kanan dan kiri) mendekati 0, sehingga $P($tidak kidal$)=1-0.09=0.91$. Dengan penalaran yang sama seperti bagian~(a), probabilitas keduanya tidak kidal adalah $0.91\times 0.91 = 0.8281$.}

\begin{exercisewrap}
\begin{nexercise}
\label{ex5Handedness}%
Misalkan 5 orang dipilih secara acak.\footnotemark\vspace{-1.5mm}
\begin{enumerate}
\setlength{\itemsep}{0mm}
\item[(a)] Berapa probabilitas semuanya tidak kidal?
\item[(b)] Berapa probabilitas semuanya kidal?
\item[(c)] Berapa probabilitas tidak semua orang tidak kidal?
\end{enumerate}
\end{nexercise}
\end{exercisewrap}
\footnotetext{(a)~Singkatan \resp{RH} dan \resp{LH} digunakan masing-masing untuk tidak kidal dan kidal. Karena semuanya independen, kita menerapkan Aturan Perkalian untuk proses independen:
\begin{align*}
P(\text{kelima orang adalah \resp{RH}})
&= P(\text{pertama = \resp{RH}, kedua = \resp{RH}, ..., kelima = \resp{RH}}) \\
&= P(\text{pertama = \resp{RH}})\times P(\text{kedua = \resp{RH}})\times  \dots \times P(\text{kelima = \resp{RH}}) \\
&= 0.91\times 0.91\times 0.91\times 0.91\times 0.91 = 0.624
\end{align*}

(b)~Dengan penalaran yang sama seperti pada~(a), $0.09\times 0.09\times 0.09\times 0.09\times 0.09 = 0.0000059$

(c)~Gunakan komplemen, $P($kelima orang \resp{RH}$)$, untuk menjawab pertanyaan ini:
\begin{align*}
P(\text{not all \resp{RH}})
	= 1 - P(\text{all \resp{RH}})
	= 1 - 0.624 = 0.376
\end{align*}}

Misalkan variabel \var{handedness} dan
\var{sex} independen,
artinya mengetahui \var{sex} seseorang tidak memberikan informasi
berguna tentang \var{handedness}-nya, dan sebaliknya.
Kita dapat menghitung probabilitas orang yang dipilih secara acak
tidak kidal dan perempuan\footnote{Proporsi aktual
  populasi AS yang \resp{perempuan} sekitar 50\%,
  sehingga kita menggunakan 0,5 sebagai probabilitas memilih perempuan.
  Namun, probabilitas ini berbeda di negara lain.}
dengan Aturan Perkalian:
\begin{align*}
P(\text{tidak kidal dan perempuan})
    &= P(\text{tidak kidal}) \times  P(\text{perempuan}) \\
    &= 0.91 \times  0.50 = 0.455
\end{align*}


\begin{exercisewrap}
\begin{nexercise}
Tiga orang dipilih secara acak.\footnotemark \vspace{-1.5mm}
\begin{enumerate}
\setlength{\itemsep}{0mm}
\item[(a)] Berapa probabilitas orang pertama laki-laki dan tidak kidal?
\item[(b)] Berapa probabilitas dua orang pertama laki-laki dan tidak kidal?
\item[(c)] Berapa probabilitas orang ketiga perempuan dan kidal?
\item[(d)] Berapa probabilitas dua orang pertama laki-laki dan tidak kidal, sedangkan orang ketiga perempuan dan kidal?
\end{enumerate}
\end{nexercise}
\end{exercisewrap}
\footnotetext{Jawaban singkat diberikan. (a)~Dalam notasi probabilitas, ini dapat ditulis $P($orang yang dipilih secara acak adalah laki-laki dan tidak kidal$)=0.455$. (b)~0.207. (c)~0.045. (d)~0.0093.}

Kadang-kadang kita ingin tahu apakah suatu hasil memberikan informasi berguna tentang hasil lain. Pertanyaannya: apakah kemunculan kedua peristiwa itu independen? Dua peristiwa $A$ dan $B$ disebut independen jika memenuhi
$P(A \text{ dan }B) = P(A) \times  P(B)$.

\begin{examplewrap}
\begin{nexample}{Jika kita mengocok dek kartu lalu mengambil satu kartu, apakah peristiwa kartu itu hati independen dari peristiwa kartu itu ace?}
Probabilitas kartu itu hati adalah $1/4$ dan probabilitas kartu itu ace adalah $1/13$. Probabilitas kartu itu ace hati adalah $1/52$.
Kita memeriksa apakah $P(A \text{ dan }B) = P(A) \times  P(B)$
terpenuhi:
\begin{align*}
P({\color{redcards}\heartsuit})\times P(\text{ace}) = \frac{1}{4}\times \frac{1}{13} = \frac{1}{52} 
					= P({\color{redcards}\heartsuit}\text{ dan ace})
\end{align*}
Karena persamaan tersebut berlaku, peristiwa kartu itu hati dan peristiwa kartu itu ace merupakan peristiwa independen.
\end{nexample}
\end{examplewrap}

% B009 translated boundary ends here: Chapter 3 opener and Defining probability
% through Independence (source lines 1--684). The separate EOC-only input below
% remains an untranslated upstream witness and is intentionally outside this
% candidate; it has no public answer payload in this source file.
{\input{ch_probability/TeX/defining_probability.tex}}




%_________________
\section{Conditional probability}
\label{conditionalProbabilitySection}

There can be rich relationships between two or more
variables that are useful to understand.
For example a car insurance company will consider
information about a person's driving history to assess
the risk that they will be responsible for an accident.
These types of relationships are the realm of conditional
probabilities.


\subsection{Exploring probabilities with a contingency table}

\index{data!photo\_classify|(}

\newcommand{\fashN}{1822}
% In order of ML, then Human
\newcommand{\fashYY}{197}
\newcommand{\fashYN}{22}
\newcommand{\fashYA}{219}
\newcommand{\fashNY}{112}
\newcommand{\fashNN}{1491}
\newcommand{\fashNA}{1603}
\newcommand{\fashAY}{309}
\newcommand{\fashAN}{1513}
\newcommand{\fashAA}{\fashN{}}
%\newcommand{\fashPYY}{}
%\newcommand{\fashPYN}{}
%\newcommand{\fashPNY}{}
%\newcommand{\fashPNN}{}
%\newcommand{\fashPYA}{0.12}
%\newcommand{\fashPNA}{0.88}
%\newcommand{\fashPAY}{}
%\newcommand{\fashPAN}{}
%\newcommand{\fashPYCY}{}
%\newcommand{\fashPYCN}{}
%\newcommand{\fashPNCY}{}
%\newcommand{\fashPNCN}{}
\newcommand{\fashCYPY}{0.96}
\newcommand{\fashCYPN}{0.04}
\newcommand{\fashCNPY}{0.07}
\newcommand{\fashCNPN}{0.93}

The \data{photo\us{}classify} data set represents
a classifier a sample of \fashN{} photos from a photo sharing website.
Data scientists have been working to improve a classifier for
whether the photo is about fashion or not, and these 1822 photos
represent a test for their classifier.
Each photo gets two classifications:
the first is called \var{mach\us{}learn} and gives
a classification from a machine
learning~(ML)\index{machine learning (ML)} system of
either \resp{pred\us{}fashion} or \resp{pred\us{}not}.
Each of these \fashN{} photos have also been classified carefully
by a team of people, which we take to be the source of truth;
this variable is called \var{truth} and takes values
\resp{fashion} and \resp{not}.
Figure~\ref{contTableOfFashionPhotos} summarizes the results.

\begin{figure}[ht]
\centering
\begin{tabular}{ll ccc rr}
&& \multicolumn{2}{c}{\var{truth}} & \hspace{1cm} &  \\
\cline{3-4}
&& \resp{fashion} & \resp{not} & Total  \\
\cline{2-5}
& \resp{pred\us{}fashion} &
    \fashYY{} & \fashYN{} & \fashYA{} \\
\raisebox{1.5ex}[0pt]{\var{mach\us{}learn}}
    & \resp{pred\us{}not} \hspace{0.5cm} &
    \fashNY{} & \fashNN{} & \fashNA{}   \\
\cline{2-5}
& Total & \fashAY{} & \fashAN{} & \fashN{} \\
\end{tabular}
\caption{Contingency table summarizing the
    \data{photo\us{}classify} data set.}
\label{contTableOfFashionPhotos}
\end{figure}
% library(openintro); table(photo_classify)

\begin{figure}[ht]
  \centering
  \Figure[A Venn diagram is shown, using boxes instead of circles, for the two categories of "ML Predicts Fashion" and "Fashion Photos" that partially overlap. The section of the rectangle for ML Predicts Fashion that is non-overlapping is labeled with 0.01. The section of the rectangle for Fashion Photos that is non-overlapping is labeled with 0.06. The overlapping section is labeled with 0.11. Outside of the rectangles is a label for "Neither" with a value 0.82.]{0.65}{photoClassifyVenn}
  \caption{A Venn diagram using boxes for the
      \data{photo\us{}classify} data set.}
  \label{photoClassifyVenn}
\end{figure}

\begin{examplewrap}
\begin{nexample}{If a photo is actually about fashion,
    what is the chance the ML classifier correctly identified
    the photo as being about fashion?}
  We can estimate this probability using the data.
  Of the \fashAY{} fashion photos,
  the ML algorithm correctly classified \fashYY{} of the photos:
\begin{align*}
P(\text{\var{mach\us{}learn} is \resp{pred\us{}fashion}
    given \var{truth} is \resp{fashion}})
  = \frac{\fashYY{}}{\fashAY{}}
  = 0.638
\end{align*}
\end{nexample}
\end{examplewrap}

\begin{examplewrap}
\begin{nexample}{We sample a photo from the data set
    and learn the ML algorithm predicted this photo
    was not about fashion.
    What is the probability that it was incorrect and
    the photo is about fashion?}
  If the ML classifier suggests a photo is not about fashion,
  then it comes from the second row in the data set.
  Of~these \fashNA{} photos, \fashNY{} were actually
  about fashion:
\begin{align*}
P(\text{\var{truth} is \resp{fashion}
    given \var{mach\us{}learn} is \resp{pred\us{}not}})
  = \frac{\fashNY{}}{\fashNA{}}
  = 0.070
\end{align*}
\end{nexample}
\end{examplewrap}

\subsection{Marginal and joint probabilities}
\label{marginalAndJointProbabilities}

\index{marginal probability|(}
\index{joint probability|(}

Figure~\ref{contTableOfFashionPhotos} includes row and
column totals for each variable separately in the
\data{photo\us{}classify} data set.
These totals represent
\termsub{marginal probabilities}{marginal probability}
for the sample, which are the probabilities based on a
single variable without regard to any other variables.
For instance, a probability based solely on the
\var{mach\us{}learn} variable is a marginal probability:
\begin{align*}
P(\text{\var{mach\us{}learn} is \resp{pred\us{}fashion}})
    = \frac{\fashYA{}}{\fashN{}}
    = 0.12
\end{align*}
A probability of outcomes for two or more variables
or processes is called a
\termsub{joint \mbox{probability}}{joint probability}:
\begin{align*}
P(\text{\var{mach\us{}learn} is \resp{pred\us{}fashion}
    and \var{truth} is \resp{fashion}})
  = \frac{\fashYY{}}{\fashN{}}
  = 0.11
\end{align*}
It is common to substitute a comma for ``and'' in a joint
probability, although using either the word ``and'' or a
comma is acceptable:
\begin{center}
$P(\text{\var{mach\us{}learn} is \resp{pred\us{}fashion},
    \var{truth} is \resp{fashion}})$ \\[2mm]
means the same thing as \\[2mm]
$P(\text{\var{mach\us{}learn} is \resp{pred\us{}fashion}
    and \var{truth} is \resp{fashion}})$
\end{center}

\begin{onebox}{Marginal and joint probabilities}
  If a probability is based on a single variable,
  it is a \emph{\hiddenterm{marginal probability}}.
  The probability of outcomes for two or more variables
  or processes is called a \emph{\hiddenterm{joint probability}}.
\end{onebox}

We use \term{table proportions} to summarize joint probabilities
for the \data{photo\us{}classify} sample.
These proportions are computed by dividing each count in
Figure~\ref{contTableOfFashionPhotos} by the table's total,
\fashN{}, to obtain the proportions in
Figure~\ref{photoClassifyProbTable}.
The joint probability distribution of the \var{mach\us{}learn}
and \var{truth} variables is shown in
Figure~\ref{photoClassifyDistribution}.

\begin{figure}[h]
\centering
\begin{tabular}{l rr r}
\hline
& \var{truth}: \resp{fashion} &
    \var{truth}: \resp{not} & Total  \\
\hline
\var{mach\us{}learn}: \resp{pred\us{}fashion} \hspace{0.5cm}
    & 0.1081 & 0.0121 & 0.1202 \\
\var{mach\us{}learn}: \resp{pred\us{}not}
    & 0.0615 & 0.8183 & 0.8798  \\
\hline
Total & 0.1696 & 0.8304 & 1.00 \\
\hline
\end{tabular}
\caption{Probability table summarizing the
    \var{photo\us{}classify} data set.}
\label{photoClassifyProbTable}
\end{figure}

\begin{figure}[h]
\centering
\begin{tabular}{l c}
  \hline
Joint outcome & Probability \\
  \hline
\var{mach\us{}learn} is \resp{pred\us{}fashion}
    and \var{truth} is \resp{fashion} & 0.1081 \\
\var{mach\us{}learn} is \resp{pred\us{}fashion}
    and \var{truth} is \resp{not} & 0.0121 \\
\var{mach\us{}learn} is \resp{pred\us{}not}
    and \var{truth} is \resp{fashion} & 0.0615 \\
\var{mach\us{}learn} is \resp{pred\us{}not}
    and \var{truth} is \resp{not} & 0.8183 \\
   \hline
Total & 1.0000 \\
\hline
\end{tabular}
\caption{Joint probability distribution for the \data{photo\us{}classify} data set.}
\label{photoClassifyDistribution}
\end{figure}

\D{\newpage}

\begin{exercisewrap}
\begin{nexercise}
Verify Figure~\ref{photoClassifyDistribution} represents
a probability distribution: events are disjoint,
all probabilities are non-negative, and the probabilities
sum to~1.\footnotemark
\end{nexercise}
\end{exercisewrap}
\footnotetext{Each of the four outcome combination are disjoint,
   all probabilities are indeed non-negative, and the sum of
   the probabilities is $0.1081 + 0.0121 + 0.0615 + 0.8183 = 1.00$.}

We can compute marginal probabilities using joint probabilities
in simple cases.
For example, the probability a randomly selected photo from the
data set is about fashion is found by summing the outcomes where
\var{truth} takes value \resp{fashion}:%
\index{marginal probability|)}\index{joint probability|)}
\newcommand{\ultruthfashion}[0]
    {\underline{\var{truth} is \resp{fashion}}}%
\begin{align*}
P(\text{\ultruthfashion{}})
  &= P(\text{\var{mach\us{}learn} is \resp{pred\us{}fashion}
        and \ultruthfashion{}}) \\
    & \qquad + P(\text{\var{mach\us{}learn} is \resp{pred\us{}not}
        and \ultruthfashion{}}) \\
  &= 0.1081 + 0.0615 \\
  &= 0.1696
\end{align*}


\subsection{Defining conditional probability}

\index{conditional probability|(}

The ML classifier predicts whether a photo is about fashion,
even if it is not perfect.
We would like to better understand how to use information
from a variable like \var{mach\us{}learn} to improve our
probability estimation of a second variable, which in this
example is \var{truth}.

The probability that a random photo from the data set is about
fashion is about 0.17.
If we knew the machine learning classifier predicted the
photo was about fashion, could we get a better estimate of the
probability the photo is actually about fashion?
Absolutely.
To do so, we limit our view to only those \fashYA{} cases
where the ML classifier predicted that the photo was about
fashion and look at the fraction where the photo was actually
about fashion:
\begin{align*}
P(\text{\var{truth} is \resp{fashion} given
    \var{mach\us{}learn} is \resp{pred\us{}fashion}})
  = \frac{\fashYY{}}{\fashYA{}}
  = 0.900
\end{align*}
We call this a \term{conditional probability} because
we computed the probability under a condition:
the ML classifier prediction said the photo was about fashion.

There are two parts to a conditional probability,
the \term{outcome of interest} and the \term{condition}.
It is useful to think of the condition as information we know
to be true, and this information usually can be described as
a known outcome or~event.
We generally separate the text inside our probability notation
into the outcome of interest and the condition with a
vertical bar:
\begin{align*}
&& P(\text{\var{truth} is \resp{fashion} given
    \var{mach\us{}learn} is \resp{pred\us{}fashion}}) \\
&& \quad = P(\text{\var{truth} is \resp{fashion}\ }|
    \text{\ \var{mach\us{}learn} is \resp{pred\us{}fashion}})
  = \frac{\fashYY{}}{\fashYA{}}
  = 0.900
\end{align*}
The vertical bar ``$|$'' is read as \emph{given}.

\D{\newpage}

In the last equation, we computed the probability a photo
was about fashion based on the condition that the ML algorithm
predicted it was about fashion as a fraction:
\begin{align*}
& P(\text{\var{truth} is \resp{fashion}\ }|
    \text{\ \var{mach\us{}learn} is \resp{pred\us{}fashion}}) \\
  &\quad = \frac{\text{\# cases where \var{truth} is \resp{fashion}
       and \var{mach\us{}learn} is \resp{pred\us{}fashion}}}
     {\text{\# cases where \var{mach\us{}learn} is \resp{pred\us{}fashion}}} \\
  &\quad = \frac{\fashYY{}}{\fashYA{}}
      = 0.900
\end{align*}
We considered only those cases that met the condition,
\var{mach\us{}learn} is \resp{pred\us{}fashion}, and then
we computed the ratio of those cases that satisfied our
outcome of interest, photo was actually about fashion.

Frequently, marginal and joint probabilities are provided
instead of count data.
For example, disease rates are commonly listed in percentages
rather than in a count format.
We would like to be able to compute conditional probabilities
even when no counts are available, and we use the last equation
as a template to understand this technique.

We considered only those cases that satisfied the condition,
where the ML algorithm predicted fashion.
Of these cases, the conditional probability was the
fraction representing the outcome of interest, that the
photo was about fashion.
Suppose we were provided only the information in
Figure~\ref{photoClassifyProbTable}, i.e. only probability data.
Then if we took a sample of 1000 photos, we would anticipate
about 12.0\% or $0.120\times 1000 = 120$ would be predicted to be
about fashion (\var{mach\us{}learn} is \resp{pred\us{}fashion}).
Similarly, we would expect about 10.8\% or
$0.108\times 1000 = 108$ to meet both the information criteria
and represent our outcome of interest.
Then the conditional probability can be computed as
\begin{align*}
&P(\text{\var{truth} is \resp{fashion}}\ |\ 
    \text{\var{mach\us{}learn} is \resp{pred\us{}fashion}}) \\
  &= \frac{\text{\# (\var{truth} is \resp{fashion}
      and \var{mach\us{}learn} is \resp{pred\us{}fashion})}}
    {\text{\# (\var{mach\us{}learn} is \resp{pred\us{}fashion})}} \\
  &= \frac{108}{120}
		= \frac{0.108}{0.120}
		= 0.90
\end{align*}
Here we are examining exactly the fraction of two probabilities,
0.108 and 0.120, which we can write as
\begin{align*}
P(\text{\var{truth} is \resp{fashion} and
    \var{mach\us{}learn} is \resp{pred\us{}fashion}})
\quad\text{and}\quad
P(\text{\var{mach\us{}learn} is \resp{pred\us{}fashion}}).
\end{align*}
The fraction of these probabilities is an example of the
general formula for conditional probability.

\begin{onebox}{Conditional probability}
  The conditional probability of outcome $A$
  given condition $B$ is computed as the following:
  \begin{align*}
  P(A | B) = \frac{P(A\text{ and }B)}{P(B)}
  \end{align*}
\end{onebox}

%\D{\newpage}

\begin{exercisewrap}
\begin{nexercise}
\label{fashionProbOfMLNotGivenTruthNot}%
(a) Write out the following statement in conditional
probability notation:
``\emph{The probability that the ML prediction was correct,
if the photo was about fashion}''.
Here the condition is now based on the photo's
\var{truth} status, not the ML algorithm. \\[1mm]
(b)~Determine the probability from part (a).
Table~\vref{photoClassifyProbTable} may be helpful.\footnotemark
\end{nexercise}
\end{exercisewrap}
\footnotetext{(a) If the photo is about fashion and the
  ML algorithm prediction was correct, then the ML algorithm
  my have a value of \resp{pred\us{}fashion}:
  \begin{align*}
  P(\text{\var{mach\us{}learn} is \resp{pred\us{}fashion}}\ |
      \ \text{\var{truth} is \resp{fashion}})
  \end{align*}
  (b)~The equation for conditional probability indicates we
  should first find \\
  $P(\text{\var{mach\us{}learn} is \resp{pred\us{}fashion}
    and \var{truth} is \resp{fashion}}) = 0.1081$
  and $P(\text{\var{truth} is \resp{fashion}}) = 0.1696$. \\
  Then the ratio represents the conditional probability:
  $0.1081 / 0.1696 = 0.6374$.}

\begin{exercisewrap}
\begin{nexercise}
\label{whyCondProbSumTo1}%
(a)~Determine the probability that the algorithm is incorrect
if it is known the photo is about fashion. \\[1mm]
(b)~Using the answers from part~(a) and
Guided Practice~\ref{fashionProbOfMLNotGivenTruthNot}(b),
compute
\begin{align*}
&P(\text{\var{mach\us{}learn} is \resp{pred\us{}fashion}}
    \ |\ \text{\var{truth} is \resp{fashion}}) \\
&\qquad  +\ P(\text{\var{mach\us{}learn} is \resp{pred\us{}not}}
    \ |\ \text{\var{truth} is \resp{fashion}})
\end{align*}
(c)~Provide an intuitive argument to explain why the sum
in~(b) is~1.\footnotemark
\end{nexercise}
\end{exercisewrap}
\footnotetext{(a)~This probability is
  $\frac{P(\text{\var{mach\us{}learn} is \resp{pred\us{}not},
      \var{truth} is \resp{fashion}})}
    {P(\text{\var{truth} is \resp{fashion}})}
  = \frac{0.0615}{0.1696} = 0.3626$.
  (b)~The total equals~1.
  (c)~Under the condition the photo is about fashion,
      the ML algorithm must have either predicted it was
      about fashion or predicted it was not about fashion.
      The complement still works for conditional probabilities,
      provided the probabilities are conditioned on the same
      information.}

\index{conditional probability|)}
\index{data!photo\_classify|)}


\subsection{Smallpox in Boston, 1721}

\index{data!smallpox|(}

The \data{smallpox} data set provides a sample of 6,224 individuals from the year 1721 who were exposed to smallpox in Boston.
Doctors at the time believed that inoculation, which involves exposing a person to the disease in a controlled form, could reduce the likelihood of death.

Each case represents one person with two variables: \var{inoculated} and \var{result}. The variable \var{inoculated} takes two levels: \resp{yes} or \resp{no}, indicating whether the person was inoculated or not. The variable \var{result} has outcomes \resp{lived} or \resp{died}. These data are summarized in Tables~\ref{smallpoxContingencyTable} and~\ref{smallpoxProbabilityTable}.

\begin{figure}[h]
\centering
\begin{tabular}{ll rr r}
& & \multicolumn{2}{c}{inoculated} & \\
\cline{3-4}
& & \resp{yes} & \resp{no} & Total  \\
\cline{2-5}
		& \resp{lived}     & 238 & 5136 & 5374 \\
\raisebox{1.5ex}[0pt]{\var{result}} &  \resp{died} \hspace{0.5cm} & 6 & 844 & 850  \\
\cline{2-5}
	& Total & 244 & 5980 & 6224 \\
\end{tabular}
\caption{Contingency table for the \data{smallpox} data set.}
\label{smallpoxContingencyTable}
\end{figure}

\begin{figure}[h]
\centering
\begin{tabular}{ll rr r}
& & \multicolumn{2}{c}{inoculated} & \\
\cline{3-4}
& & \resp{yes} & \resp{no} & Total  \\
   \cline{2-5}
 & \resp{lived}     & 0.0382 & 0.8252 & 0.8634 \\
\raisebox{1.5ex}[0pt]{\var{result}} & \resp{died} \hspace{0.5cm} & 0.0010 & 0.1356  & 0.1366  \\
   \cline{2-5}
& Total & 0.0392 & 0.9608 & 1.0000 \\
\end{tabular}
\caption{Table proportions for the \data{smallpox} data, computed by dividing each count by the table total, 6224.}
\label{smallpoxProbabilityTable}
\end{figure}

%\D{\newpage}

\begin{exercisewrap}
\begin{nexercise} \label{probDiedIfNotInoculated}
Write out, in formal notation, the probability a randomly selected person who was not inoculated died from smallpox, and find this \mbox{probability.}\footnotemark
\end{nexercise}
\end{exercisewrap}
\footnotetext{$P($\var{result} = \resp{died} $|$ \var{inoculated} = \resp{no}$) = \frac{P(\text{\var{result} = \resp{died} and \var{inoculated} = \resp{no}})}{P(\text{\var{inoculated} = \resp{no}})} = \frac{0.1356}{0.9608} = 0.1411$.}

\begin{exercisewrap}
\begin{nexercise}
Determine the probability that an inoculated person died from smallpox. How does this result compare with the result of Guided Practice~\ref{probDiedIfNotInoculated}?\footnotemark
\end{nexercise}
\end{exercisewrap}
\footnotetext{$P($\var{result} = \resp{died} $|$ \var{inoculated} = \resp{yes}$) = \frac{P(\text{\var{result} = \resp{died} and \var{inoculated} = \resp{yes}})}{P(\text{\var{inoculated} = \resp{yes}})} = \frac{0.0010}{0.0392} = 0.0255$ (if we avoided rounding errors, we'd get $6 / 244 = 0.0246$). The death rate for individuals who were inoculated is only about 1~in~40 while the death rate is about 1~in~7 for those who were not inoculated.}

\begin{exercisewrap}
\begin{nexercise}\label{SmallpoxInoculationObsExpExercise}
The people of Boston self-selected whether or not to be inoculated. (a) Is this study observational or was this an experiment? (b) Can we infer any causal connection using these data? (c) What are some potential confounding variables that might influence whether someone \resp{lived} or \resp{died} and also affect whether that person was inoculated?\footnotemark
\end{nexercise}
\end{exercisewrap}
\footnotetext{Brief answers: (a)~Observational. (b)~No, we cannot infer causation from this observational study. (c)~Accessibility to the latest and best medical care. There are other valid answers for part~(c).}


\subsection{General multiplication rule}

Section~\ref{probabilityIndependence} introduced the Multiplication Rule for independent processes. Here we provide the \term{General Multiplication Rule} for events that might not be independent.

\begin{onebox}{General Multiplication Rule}
If $A$ and $B$ represent two outcomes or events, then \vspace{-1.5mm}
\begin{align*}
P(A\text{ and }B) = P(A | B)\times P(B)
\end{align*} \vspace{-6.5mm} \par
It is useful to think of $A$ as the outcome of interest and $B$ as the condition.
\end{onebox}

\noindent%
This General Multiplication Rule is simply a rearrangement of the conditional probability equation.

%\D{\newpage}

\begin{examplewrap}
\begin{nexample}{Consider the \data{smallpox} data set. Suppose we are given only two pieces of information: 96.08\% of residents were not inoculated, and 85.88\% of the residents who were not inoculated ended up surviving. How could we compute the probability that a resident was not inoculated and lived?}
We will compute our answer using the General Multiplication Rule and then verify it using Figure~\ref{smallpoxProbabilityTable}. We want to determine
\begin{align*}
P(\text{\var{result}
    = \resp{lived} and \var{inoculated} = \resp{no}})
\end{align*}
and we are given that
\begin{align*}
P(\text{\var{result}
    = \resp{lived} }|\text{ \var{inoculated} = \resp{no}})
    &= 0.8588 %\\
&&P(\text{\var{inoculated} = \resp{no}})
    = 0.9608
\end{align*}
Among the 96.08\% of people who were not inoculated, 85.88\% survived:
\begin{align*}
P(\text{\var{result} = \resp{lived}
        and \var{inoculated} = \resp{no}})
    = 0.8588 \times 0.9608
    = 0.8251
\end{align*}
This is equivalent to the General Multiplication Rule. We can confirm this probability in Figure~\ref{smallpoxProbabilityTable} at the intersection of \resp{no} and \resp{lived} (with a small rounding error).
\end{nexample}
\end{examplewrap}

\begin{exercisewrap}
\begin{nexercise}
Use $P($\var{inoculated} = \resp{yes}$) = 0.0392$ and $P($\var{result} = \resp{lived} $|$ \var{inoculated} = \resp{yes}$) = 0.9754$ to determine the probability that a person was both inoculated and lived.\footnotemark
\end{nexercise}
\end{exercisewrap}
\footnotetext{The answer is 0.0382, which can be verified using Figure~\ref{smallpoxProbabilityTable}.}

%\D{\newpage}

\begin{exercisewrap}
\begin{nexercise}
If 97.54\% of the inoculated people lived,
what proportion of inoculated people must have died?\footnotemark{}
\end{nexercise}
\end{exercisewrap}
\footnotetext{There were only two possible outcomes: \resp{lived} or \resp{died}. This means that 100\% - 97.54\% = 2.46\% of the people who were inoculated died.}

\begin{onebox}{Sum of conditional probabilities}
Let $A_1$, ..., $A_k$ represent all the disjoint outcomes for a variable or process. Then if $B$ is an event, possibly for another variable or process, we have: \vspace{-1mm}
\begin{align*}
P(A_1|B) + \cdots + P(A_k|B) = 1
\end{align*}%
\vspace{-5.5mm} \par
The rule for complements also holds when an event and its complement are conditioned on the same information: \vspace{-1.5mm}
\begin{align*}
P(A | B) = 1 - P(A^c | B)
\end{align*}
\end{onebox}

\begin{exercisewrap}
\begin{nexercise}
Based on the probabilities computed above, does it appear that inoculation is effective at reducing the risk of death from smallpox?\footnotemark
\end{nexercise}
\end{exercisewrap}
\footnotetext{The samples are large relative to the difference in death rates for the ``inoculated'' and ``not inoculated'' groups, so it seems there is an association between \var{inoculated} and \var{outcome}. However, as noted in the solution to Guided Practice~\ref{SmallpoxInoculationObsExpExercise}, this is an observational study and we cannot be sure if there is a causal connection. (Further research has shown that inoculation is effective at reducing death rates.)}


%\D{\newpage}

\subsection{Independence considerations in conditional probability}

If two events are independent, then knowing the outcome of one should provide no information about the other. We can show this is mathematically true using conditional probabilities.

\begin{exercisewrap}
\begin{nexercise} \label{condProbOfRollingA1AfterOne1}
Let $X$ and $Y$ represent the outcomes of rolling two dice.\footnotemark
\begin{enumerate}[(a)]
\setlength{\itemsep}{0mm}
\item What is the probability that the first die, $X$, is \resp{1}?
\item What is the probability that both $X$ and $Y$ are \resp{1}?
\item Use the formula for conditional probability to compute $P(Y =$ \resp{1}$\ |\ X = $ \resp{1}$)$.
\item What is $P(Y=1)$? Is this different from the answer from part (c)? Explain.
\end{enumerate}
\end{nexercise}
\end{exercisewrap}
\footnotetext{Brief solutions: (a) $1/6$. (b) $1/36$. (c)~$\frac{P(Y = \text{ \resp{1} and }X=\text{ \resp{1}})}{P(X=\text{ \resp{1}})} = \frac{1/36}{1/6} = 1/6$. (d)~The probability is the same as in part~(c): $P(Y=1)=1/6$. The probability that $Y=1$ was unchanged by knowledge about $X$, which makes sense as $X$ and $Y$ are independent.}

We can show in Guided Practice~\ref{condProbOfRollingA1AfterOne1}(c) that the conditioning information has no influence by using the Multiplication Rule for independence processes:
\begin{align*}
P(Y=\text{\resp{1}}\ |\ X=\text{\resp{1}})
    &= \frac{P(Y=\text{\resp{1} and }X=\text{\resp{1}})}
      {P(X=\text{\resp{1}})} \\
    &= \frac{P(Y=\text{\resp{1}}) \times
        \color{oiGB}P(X=\text{\resp{1}})}
      {\color{oiGB}P(X=\text{\resp{1}})} \\
    &= P(Y=\text{\resp{1}}) \\
\end{align*}

\begin{exercisewrap}
\begin{nexercise}
Ron is watching a roulette table in a casino and notices that the last five outcomes were \resp{black}. He figures that the chances of getting \resp{black} six times in a row is very small (about $1/64$) and puts his paycheck on red. What is wrong with his reasoning?\footnotemark
\end{nexercise}
\end{exercisewrap}
\footnotetext{He has forgotten that the next roulette spin is independent of the previous spins. Casinos do employ this practice, posting the last several outcomes of many betting games to trick unsuspecting gamblers into believing the odds are in their favor. This is called the \term{gambler's fallacy}.}


\D{\newpage}

\subsection{Tree diagrams}

\index{data!smallpox|)}
\index{tree diagram|(}

\termsub{Tree diagrams}{tree diagram} are a tool to organize outcomes and probabilities around the structure of the data. They are most useful when two or more processes occur in a sequence and each process is conditioned on its predecessors.

The \data{smallpox} data fit this description. We see the population as split by \var{inoculation}: \resp{yes} and \resp{no}. Following this split, survival rates were observed for each group. This structure is reflected in the \term{tree diagram} shown in Figure~\ref{smallpoxTreeDiagram}. The first branch for \var{inoculation} is said to be the \term{primary} branch while the other branches are \termni{secondary}.

\begin{figure}[ht]
\centering
\Figure[A tree diagram with a primary branch "Inoculated" and a secondary branch "Result". The Inoculated primary branching leads to two options: "Yes" with a probability of 0.0392 and "No" with a probability of 0.9608. Each of these branches has secondary branches with conditional probabilities for the "Result" conditional on "Inoculated". The Inoculated Yes branch breaks into branches for "Lived" (0.9754) and "Died" (0.0246). These branches also provide the multiplied probabilities along the branches as well. For example, the Yes-and-Lived branching multiplies 0.0392 times 0.9754 to get 0.03824. The Yes-and-Died branching has a multiplied probability of 0.00096. Next, turning our attention to the "No" primary branch, it also has secondary branches of Lived and Died with conditional probabilities 0.8589 and 0.1411, respectively. It also shows the probabilities multiplied along each set of branches, with No-and-Lived as 0.82523 and No-and-Died as 0.13557.]{0.93}{smallpoxTreeDiagram}
\caption{A tree diagram of the \data{smallpox} data set.}
\label{smallpoxTreeDiagram}
\end{figure}

Tree diagrams are annotated with marginal and conditional probabilities, as shown in Figure~\ref{smallpoxTreeDiagram}. This tree diagram splits the smallpox data by \var{inoculation} into the \resp{yes} and \resp{no} groups with respective marginal probabilities 0.0392 and 0.9608. The secondary branches are conditioned on the first, so we assign conditional probabilities to these branches. For example, the top branch in Figure~\ref{smallpoxTreeDiagram} is the probability that \var{result} = \resp{lived} conditioned on the information that \var{inoculated} = \resp{yes}. We may (and usually do) construct joint probabilities at the end of each branch in our tree by multiplying the numbers we come across as we move from left to right. These joint probabilities are computed using the General Multiplication Rule:
\begin{align*}
& P(\text{\var{inoculated} = \resp{yes}
    and \var{result} = \resp{lived}}) \\
  &\quad = P(\text{\var{inoculated} = \resp{yes}})\times
      P(\text{\var{result} = \resp{lived}}|
          \text{\var{inoculated} = \resp{yes}}) \\
  &\quad = 0.0392\times 0.9754=0.0382
\end{align*}

\begin{examplewrap}
\begin{nexample}{Consider the midterm and final for a statistics class. Suppose 13\% of students earned an \resp{A} on the midterm. Of those students who earned an \resp{A} on the midterm, 47\% received an \resp{A} on the final, and 11\% of the students who earned lower than an \resp{A} on the midterm received an \resp{A} on the final. You randomly pick up a final exam and notice the student received an \resp{A}. What is the probability that this student earned an \resp{A} on the midterm?} \label{exerciseForTreeDiagramOfStudentGettingAOnMidtermGivenThatSheGotAOnFinal}
The end-goal is to find $P(\text{\var{midterm} = \resp{A}} | \text{\var{final} = \resp{A}})$. To calculate this conditional probability, we need the following probabilities:
\begin{align*}
P(\text{\var{midterm} = \resp{A} and \var{final} = \resp{A}})
  \qquad\text{and}\qquad
  P(\text{\var{final} = \resp{A}})
\end{align*}
However, this information is not provided, and it is not obvious how to calculate these probabilities. Since we aren't sure how to proceed, it is useful to organize the information into a tree diagram:
\begin{center}
\Figure[A tree diagram with a primary branch "Midterm" and a secondary branch "Final". The Midterm primary branching leads to two options: "A" with a probability of 0.13 and "Other" with a probability of 0.87. Each of these branches has secondary branches with conditional probabilities for the "Final" conditional on "Midterm". The Midterm-A branch breaks into branches for "A", with a conditional probability of 0.47 with an A-and-A final probability of 0.0611, and an "other" secondary branch, with a conditional probability of 0.53 with an other-and-other final probability of 0.0689. Next, turning our attention to the Midterm-Other primary branch, it also has secondary branches of Final-A with a conditional probability of 0.11 and final probability of 0.0957, and an "Final-other" branch with a conditional probability of 0.89 and final probability of 0.7743.]{0.85}{testTree}
\end{center}
When constructing a tree diagram, variables provided with marginal probabilities are often used to create the tree's primary branches; in this case, the marginal probabilities are provided for midterm grades. The final grades, which correspond to the conditional probabilities provided, will be shown on the secondary branches.

With the tree diagram constructed, we may compute the required probabilities:
\begin{align*}
&P(\text{\var{midterm} = \resp{A} and \var{final} = \resp{A}}) = 0.0611 \\
&P(\text{\underline{\color{black}\var{final} = \resp{A}}})  \\
& \quad= P(\text{\var{midterm} = \resp{other} and \underline{\color{black}\var{final} = \resp{A}}}) + P(\text{\var{midterm} = \resp{A} and \underline{\color{black}\var{final} = \resp{A}}}) \\
& \quad= 0.0957 + 0.0611  = 0.1568
\end{align*}
The marginal probability, $P($\var{final} = \resp{A}$)$, was calculated by adding up all the joint probabilities on the right side of the tree that correspond to \var{final} = \resp{A}. We may now finally take the ratio of the two probabilities:
\begin{align*}
P(\text{\var{midterm} = \resp{A}} | \text{\var{final} = \resp{A}})
  &= \frac{P(\text{\var{midterm} = \resp{A}
      and \var{final} = \resp{A}})}
    {P(\text{\var{final} = \resp{A}})} \\
  &= \frac{0.0611}{0.1568} = 0.3897
\end{align*}
The probability the student also earned an A on the midterm is about 0.39.
\end{nexample}
\end{examplewrap}

%\begin{figure}[ht]
%  \centering
%  \Figure{0.9}{testTree}
%  \caption{A tree diagram describing the \var{midterm}
%      and \var{final} variables.}
%  \label{testTree}
%\end{figure}

\begin{exercisewrap}
\begin{nexercise}
After an introductory statistics course, 78\% of students can successfully construct tree diagrams. Of those who can construct tree diagrams, 97\% passed, while only 57\% of those students who could not construct tree diagrams passed. (a)~Organize this information into a tree diagram. (b)~What is the probability that a randomly selected student passed? (c)~Compute the probability a student is able to construct a tree diagram if it is known that she passed.\footnotemark
\end{nexercise}
\end{exercisewrap}
\footnotetext{%\begin{minipage}[t]{0.27\linewidth}
(a) The tree diagram is shown to the right. \\
(b)~Identify which two joint probabilities represent
    students who passed, and add them:
    $P($passed$) = 0.7566+0.1254= 0.8820$. \\
(c)~$P($construct tree diagram $|$ passed$)
   = \frac{0.7566}{0.8820} = 0.8578$. \\ %\vspace{15mm} \\
%\end{minipage}
%\begin{minipage}[c]{0.7\linewidth}
\Figure[A tree diagram with a primary branch "Able to construct tree diagrams" and a secondary branch "Pass class". The Able-to-construct-tree-diagrams primary branching leads to two options: "Yes" with a probability of 0.78 and "No" with a probability of 0.22. Each of these branches has secondary branches with conditional probabilities for "Pass Class" conditional on "Able to construct tree diagrams". The Yes primary branch breaks into branches for "Pass", with a conditional probability of 0.97 with a Yes-and-Pass final probability of 0.7566, and a "Fail" secondary branch, with a conditional probability of 0.03 with a Yes-and-Fail final probability of 0.0234. Next, turning our attention to the No primary branch, it also has secondary branches of Pass with a conditional probability of 0.57 and final probability of 0.1254, and a Fail branch with a conditional probability of 0.43 and final probability of 0.0946.]{0.7}{treeDiagramAndPass}}% \vspace{-25mm}
%\end{minipage}}


\subsection{Bayes' Theorem}
\label{bayesTheoremSubsection}

\index{Bayes' Theorem|(}

In many instances, we are given a conditional probability of the form
\begin{align*}
P(\text{statement about variable 1 } | \text{ statement about variable 2})
\end{align*}
but we would really like to know the inverted conditional probability:
\begin{align*}
P(\text{statement about variable 2 } | \text{ statement about variable 1})
\end{align*}
Tree diagrams can be used to find the second conditional probability when given the first. However, sometimes it is not possible to draw the scenario in a tree diagram. In these cases, we can apply a very useful and general formula: Bayes' Theorem.

We first take a critical look at an example of inverting conditional probabilities where we still apply a tree diagram.

\D{\newpage}

\begin{examplewrap}
\begin{nexample}{In Canada, about 0.35\% of women over 40
    will develop breast cancer in any given year.
    A common screening test for cancer is the mammogram,
    but this test is not perfect.
    In about 11\% of patients with breast cancer, the test
    gives a \term{false negative}:
    it indicates a woman does not have breast cancer when
    she does have breast cancer.
    Similarly, the test gives a \term{false positive}
    in 7\% of patients who do not have breast cancer:
    it indicates these patients have breast cancer when
    they actually do not.
    If we tested a random woman over 40 for breast cancer
    using a mammogram and the test came back positive
    -- that is, the test suggested the patient has cancer --
    what is the probability that the patient actually has
    breast cancer?}

\label{probBreastCancerGivenPositiveTestExample}

Notice that we are given sufficient information to quickly compute the probability of testing positive if a woman has breast cancer ($1.00-0.11=0.89$). However, we seek the inverted probability of cancer given a positive test result. (Watch out for the non-intuitive medical language: a~\emph{positive} test result suggests the possible presence of cancer in a mammogram screening.) This inverted probability may be broken into two pieces:
\begin{align*}
P(\text{has BC } | \text{ mammogram$^+$}) = \frac{P(\text{has BC and mammogram$^+$})}{P(\text{mammogram$^+$})}
\end{align*}
where ``has BC'' is an abbreviation for the patient having
breast cancer and ``mammogram$^+$'' means the mammogram screening
was positive.
We can construct a tree diagram for these probabilities:
\begin{center}
\Figure[A tree diagram with a primary branch "Truth" and a secondary branch "Mammogram". The Truth primary branching leads to two options: "Cancer" with a probability of 0.0035 and "No Cancer" with a probability of 0.9965. Each of these branches has secondary branches with conditional probabilities for "Positive" and "Negative" mammogram outcomes conditional on whether the truth is having cancer or not. The Cancer primary branch breaks into branches for "Positive", with a conditional probability of 0.89 with a Cancer-and-Positive final probability of 0.00312, and a "Negative" secondary branch, with a conditional probability of 0.11 with a Cancer-and-Negative final probability of 0.00038. Next, turning our attention to the No-Cancer primary branch, it also has secondary branches of Positive with a conditional probability of 0.07 and final probability of 0.06976, and a Negative branch with a conditional probability of 0.93 and final probability of 0.92675.]{0.9}{BreastCancerTreeDiagram}
\end{center}
The probability the patient has breast cancer
and the mammogram is positive is
\begin{align*}
P(\text{has BC and mammogram$^+$}) &= P(\text{mammogram$^+$ } | \text{ has BC})P(\text{has BC}) \\
	&= 0.89\times 0.0035 = 0.00312
\end{align*}
The probability of a positive test result is the sum of the two corresponding scenarios:
\begin{align*}
P(\text{\underline{\color{black}mammogram$^+$}})
  &= P(\text{\underline{\color{black}mammogram$^+$} and has BC}) \\
  &\qquad\qquad
      + P(\text{\underline{\color{black}mammogram$^+$} and no BC})\\
  &= P(\text{has BC})P(\text{mammogram$^+$ } | \text{ has BC}) \\
  &\qquad\qquad
      + P(\text{no BC})P(\text{mammogram$^+$ } | \text{ no BC}) \\
  &= 0.0035\times 0.89 + 0.9965\times 0.07 = 0.07288
\end{align*}
Then if the mammogram screening is positive for a patient, the probability the patient has breast cancer is
\begin{align*}
P(\text{has BC } | \text{ mammogram$^+$})
	&= \frac{P(\text{has BC and mammogram$^+$})}{P(\text{mammogram$^+$})}\\
	&= \frac{0.00312}{0.07288} \approx 0.0428
\end{align*}
That is, even if a patient has a positive mammogram screening, there is still only a~4\%~chance that she has breast cancer.
\end{nexample}
\end{examplewrap}

%\begin{figure}[h]
%  \centering
%  \Figure{0.75}{BreastCancerTreeDiagram}
%  \caption{Tree diagram for
%      Example~\ref{probBreastCancerGivenPositiveTestExample}.}%,
%      %computing the probability a random patient who tests
%      %positive on a mammogram actually has breast cancer.}
%\label{BreastCancerTreeDiagram}
%\end{figure}

\D{\newpage}

Example~\ref{probBreastCancerGivenPositiveTestExample} highlights why doctors often run more tests regardless of a first positive test result. When a medical condition is rare, a single positive test isn't generally definitive.

Consider again the last equation of Example~\ref{probBreastCancerGivenPositiveTestExample}.
Using the tree diagram, we can see that the numerator (the top of the fraction) is equal to the following product:
\begin{align*}
P(\text{has BC and mammogram$^+$}) = P(\text{mammogram$^+$ } | \text{ has BC})P(\text{has BC})
\end{align*}
The denominator -- the probability the screening was positive -- is equal to the sum of probabilities for each positive screening scenario:
\begin{align*}
P(\text{\underline{\color{black}mammogram$^+$}})
	&= P(\text{\underline{\color{black}mammogram$^+$} and no BC})
		+ P(\text{\underline{\color{black}mammogram$^+$} and has BC})
\end{align*}
In the example, each of the probabilities on the right side was broken down into a product of a conditional probability and marginal probability using the tree diagram.
\begin{align*}
P(\text{mammogram$^+$})
	&= P(\text{mammogram$^+$ and no BC}) + P(\text{mammogram$^+$ and has BC}) \\
	&= P(\text{mammogram$^+$ } | \text{ no BC})P(\text{no BC}) \\
			   &\qquad\qquad + P(\text{mammogram$^+$ } | \text{ has BC})P(\text{has BC})
\end{align*}
We can see an application of Bayes' Theorem by substituting the resulting probability expressions into the numerator and denominator of the original conditional probability.
\begin{align*}
& P(\text{has BC } | \text{ mammogram$^+$})  \\
& \qquad= \frac{P(\text{mammogram$^+$ } | \text{ has BC})P(\text{has BC})}
	{P(\text{mammogram$^+$ } | \text{ no BC})P(\text{no BC}) + P(\text{mammogram$^+$ } | \text{ has BC})P(\text{has BC})}
\end{align*}

\begin{onebox}{Bayes' Theorem: inverting probabilities}
Consider the following conditional probability for variable 1 and variable 2:\vspace{-1.5mm}
\begin{align*}
P(\text{outcome $A_1$ of variable 1 } | \text{ outcome $B$ of variable 2})
\end{align*}
Bayes' Theorem states that this conditional probability can be identified as the following fraction:\vspace{-1.5mm}
\begin{align*}
\frac{P(B | A_1) P(A_1)}
	{P(B | A_1) P(A_1) + P(B | A_2) P(A_2) + \cdots + P(B | A_k) P(A_k)}
\end{align*}
where $A_2$, $A_3$, ..., and $A_k$ represent all other possible outcomes of the first variable.\index{Bayes' Theorem|textbf}
\end{onebox}

Bayes' Theorem is a generalization of what we have done
using tree diagrams.
The numerator identifies the probability of getting both
$A_1$ and~$B$.
The denominator is the marginal probability of getting~$B$.
This bottom component of the fraction appears long and
complicated since we have to add up probabilities from all of the different ways to get $B$. We always completed this step when using tree diagrams. However, we usually did it in a separate step so it didn't seem as complex.

\noindent%
To apply Bayes' Theorem correctly, there are two preparatory steps:
\begin{enumerate}
\setlength{\itemsep}{0mm}
\item[(1)] First identify the marginal probabilities of each possible outcome of the first variable: $P(A_1)$, $P(A_2)$, ..., $P(A_k)$.
\item[(2)] Then identify the probability of the outcome $B$, conditioned on each possible scenario for the first variable: $P(B | A_1)$, $P(B | A_2)$, ..., $P(B | A_k)$.
\end{enumerate}
Once each of these probabilities are identified, they can be applied directly within the formula.
Bayes' Theorem tends to be a good option when there are so many scenarios that drawing a tree diagram would be complex.

\begin{exercisewrap}
\begin{nexercise} \label{exerciseForParkingLotOnCampusBeingFullAndWhetherOrNotThereIsASportingEvent}
Jose visits campus every Thursday evening. However, some days the parking garage is full, often due to college events. There are academic events on 35\% of evenings, sporting events on 20\% of evenings, and no events on 45\% of evenings. When there is an academic event, the garage fills up about 25\% of the time, and it fills up 70\% of evenings with sporting events. On evenings when there are no events, it only fills up about 5\% of the time. If Jose comes to campus and finds the garage full, what is the probability that there is a sporting event? Use a tree diagram to solve this problem.\footnotemark
\end{nexercise}
\end{exercisewrap}
\footnotetext{\begin{minipage}[t]{0.27\linewidth}
The tree diagram, with three primary branches, is shown to the right. Next, we identify two probabilities from the tree diagram. (1) The probability that there is a sporting event and the garage is full: 0.14. (2) The probability the garage is full: $0.0875 + 0.14 + 0.0225 = 0.25$. Then the solution is the ratio of these probabilities: $\frac{0.14}{0.25} = 0.56$. If the garage is full, there is a 56\% probability that there is a sporting event. \vspace{0.1mm} \\\ 
\end{minipage}
\begin{minipage}[c]{0.65\linewidth}
\Figure[A tree diagram with a primary branch "Event" and a secondary branch "Garage full". The primary "Event" branching has three possibilities of "Academic" with probability 0.35, "Sporting" with probability 0.20, and "None" with probability 0.45. Each of these three branches has two secondary branches. The "Academic" primary branch breaks into branches for "Full" that has a conditional probability of 0.25 with an Academic-and-Full final probability of 0.0875, and a "Spaces Available" secondary branch with a conditional probability of 0.75 with an Academic-and-Spaces-Available final probability of 0.2625. The "Sporting" primary branch breaks into branches for "Full" that has a conditional probability of 0.7 with a Sporting-and-Full final probability of 0.14, and a "Spaces Available" secondary branch with a conditional probability of 0.3 with a Sporting-and-Spaces-Available final probability of 0.06. The "None" primary branch breaks into branches for "Full" that has a conditional probability of 0.05 with a None-and-Full final probability of 0.0225, and a "Spaces Available" secondary branch with a conditional probability of 0.95 with a None-and-Spaces-Available final probability of 0.4275.]{}{treeDiagramGarage}\vspace{-45mm}
\end{minipage}}

\begin{examplewrap}
\begin{nexample}{Here we solve the same problem presented in Guided Practice~\ref{exerciseForParkingLotOnCampusBeingFullAndWhetherOrNotThereIsASportingEvent}, except this time we use Bayes' Theorem.}
The outcome of interest is whether there is a sporting event (call this $A_1$), and the condition is that the lot is full ($B$). Let $A_2$ represent an academic event and $A_3$ represent there being no event on campus. Then the given probabilities can be written as
\begin{align*}
&P(A_1) = 0.2 &&P(A_2) = 0.35 &&P(A_3) = 0.45 \\
&P(B | A_1) = 0.7 &&P(B | A_2) = 0.25 &&P(B | A_3) = 0.05
\end{align*}
Bayes' Theorem can be used to compute the probability of a sporting event ($A_1$) under the condition that the parking lot is full ($B$):
\begin{align*}
P(A_1 | B) &= \frac{P(B | A_1) P(A_1)}{P(B | A_1) P(A_1) + P(B | A_2) P(A_2) + P(B | A_3) P(A_3)} \\
		&= \frac{(0.7)(0.2)}{(0.7)(0.2) + (0.25)(0.35) + (0.05)(0.45)} \\
		&= 0.56 
\end{align*}
Based on the information that the garage is full, there is a 56\% probability that a sporting event is being held on campus that evening.
\end{nexample}
\end{examplewrap}

\D{\newpage}

\begin{exercisewrap}
\begin{nexercise} \label{exerciseForParkingLotOnCampusBeingFullAndWhetherOrNotThereIsAnAcademicEvent}
Use the information in the previous exercise and example to verify the probability that there is an academic event conditioned on the parking lot being full is 0.35.\footnotemark
\end{nexercise}
\end{exercisewrap}
\footnotetext{Short answer:
\begin{align*}
P(A_2 | B) &= \frac{P(B | A_2) P(A_2)}{P(B | A_1) P(A_1) + P(B | A_2) P(A_2) + P(B | A_3) P(A_3)} \\
		&= \frac{(0.25)(0.35)}{(0.7)(0.2) + (0.25)(0.35) + (0.05)(0.45)} \\
		&= 0.35
\end{align*}}

\begin{exercisewrap}
\begin{nexercise} \label{exerciseForParkingLotOnCampusBeingFullAndWhetherOrNotThereIsNoEvent}
In Guided Practice~\ref{exerciseForParkingLotOnCampusBeingFullAndWhetherOrNotThereIsASportingEvent} and~\ref{exerciseForParkingLotOnCampusBeingFullAndWhetherOrNotThereIsAnAcademicEvent}, you found that if the parking lot is full, the probability there is a sporting event is 0.56 and the probability there is an academic event is 0.35. Using this information, compute $P($no event $|$ the lot is full$)$.\footnotemark
\end{nexercise}
\end{exercisewrap}
\footnotetext{Each probability is conditioned on the same information that the garage is full, so the complement may be used: $1.00 - 0.56 - 0.35 = 0.09$.}

The last several exercises offered a way to update our belief about whether there is a sporting event, academic event, or no event going on at the school based on the information that the parking lot was full. This strategy of \emph{updating beliefs} using Bayes' Theorem is actually the foundation of an entire section of statistics called \term{Bayesian statistics}. While Bayesian statistics is very important and useful, we will not have time to cover much more of it in this book.

\index{Bayes' Theorem|)}
\index{tree diagram|)}
\index{conditional probability|)}
\index{probability|)}


{\input{ch_probability/TeX/conditional_probability.tex}}





%_________________
\section{Sampling from a small population}
\label{smallPop}

\noindent%
When we sample observations from a population,
usually we're only sampling a small fraction of
the possible individuals or cases.
However, sometimes our sample size is large enough
or the population is small enough that we
sample more than 10\% of a population\footnote{The 10\%
  guideline is a rule of thumb cutoff for when these
  considerations become more important.}
\emph{without
replacement} (meaning we do not have a chance of
sampling the same cases twice).
Sampling such a notable fraction of a population
can be important for how we analyze the
sample.

\begin{examplewrap}
\begin{nexample}{Professors sometimes select a student at random to answer a question. If each student has an equal chance of being selected and there are 15 people in your class, what is the chance that she will pick you for the next question?}
If there are 15 people to ask and none are skipping class, then the probability is $1/15$, or about $0.067$.
\end{nexample}
\end{examplewrap}

\begin{examplewrap}
\begin{nexample}{If the professor asks 3 questions, what is the probability that you will not be selected? Assume that she will not pick the same person twice in a given lecture.}\label{3woRep}
For the first question, she will pick someone else with probability $14/15$. When she asks the second question, she only has 14 people who have not yet been asked. Thus, if you were not picked on the first question, the probability you are again not picked is $13/14$. Similarly, the probability you are again not picked on the third question is $12/13$, and the probability of not being picked for any of the three questions is
\begin{align*}
&P(\text{not picked in 3 questions}) \\
&\quad = P(\text{\var{Q1}} = \text{\resp{not\us{}picked}, }\text{\var{Q2}} = \text{\resp{not\us{}picked}, }\text{\var{Q3}} = \text{\resp{not\us{}picked}.}) \\
&\quad = \frac{14}{15}\times\frac{13}{14}\times\frac{12}{13} = \frac{12}{15} = 0.80
\end{align*}
\end{nexample}
\end{examplewrap}

\begin{exercisewrap}
\begin{nexercise}
What rule permitted us to multiply the probabilities in Example~\ref{3woRep}?\footnotemark
\end{nexercise}
\end{exercisewrap}
\footnotetext{The three probabilities we computed were actually one marginal probability, $P($\var{Q1}$ = $\resp{not\us{}picked}$)$, and two conditional probabilities:
\begin{align*}
&P(\text{\var{Q2}} =  \text{\resp{not\us{}picked} }|
    \text{ \var{Q1}} = \text{\resp{not\us{}picked}}) \\
&P(\text{\var{Q3}} =  \text{\resp{not\us{}picked} }|
    \text{ \var{Q1}} = \text{\resp{not\us{}picked}, }
    \text{\var{Q2}} = \text{\resp{not\us{}picked}})
\end{align*}
Using the General Multiplication Rule, the product of these three probabilities is the probability of not being picked in 3 questions.}

\D{\newpage}

\begin{examplewrap}
\begin{nexample}{Suppose the professor randomly picks without regard to who she already selected, i.e. students can be picked more than once. What is the probability that you will not be picked for any of the three questions?}\label{3wRep}
Each pick is independent, and the probability of not being picked for any individual question is $14/15$. Thus, we can use the Multiplication Rule for independent processes.
\begin{align*}
&P(\text{not picked in 3 questions}) \\
&\quad = P(\text{\var{Q1}} = \text{\resp{not\us{}picked}, }\text{\var{Q2}} = \text{\resp{not\us{}picked}, }\text{\var{Q3}} = \text{\resp{not\us{}picked}.}) \\
&\quad = \frac{14}{15}\times\frac{14}{15}\times\frac{14}{15} = 0.813
\end{align*}
You have a slightly higher chance of not being picked compared to when she picked a new person for each question. However, you now may be picked more than once.
\end{nexample}
\end{examplewrap}

\begin{exercisewrap}
\begin{nexercise}
Under the setup of Example~\ref{3wRep}, what is the probability of being picked to answer all three questions?\footnotemark
\end{nexercise}
\end{exercisewrap}
\footnotetext{$P($being picked to answer all three questions$) = \left(\frac{1}{15}\right)^3 = 0.00030$.}

If we sample from a small population \term{without replacement}, we no longer have independence between our observations. In Example~\ref{3woRep}, the probability of not being picked for the second question was conditioned on the event that you were not picked for the first question. In Example~\ref{3wRep}, the professor sampled her students \term{with replacement}: she repeatedly sampled the entire class without regard to who she already picked. 

\begin{exercisewrap}
\begin{nexercise} \label{raffleOf30TicketsWWOReplacement}
Your department is holding a raffle. They sell 30 tickets and offer seven prizes. (a) They place the tickets in a hat and draw one for each prize. The tickets are sampled without replacement, i.e. the selected tickets are not placed back in the hat. What is the probability of winning a prize if you buy one ticket? (b)~What if the tickets are sampled with replacement?\footnotemark
\end{nexercise}
\end{exercisewrap}
\footnotetext{(a) First determine the probability of not winning. The tickets are sampled without replacement, which means the probability you do not win on the first draw is $29/30$, $28/29$ for the second, ..., and $23/24$ for the seventh. The probability you win no prize is the product of these separate probabilities: $23/30$. That is, the probability of winning a prize is $1 - 23/30 = 7/30 = 0.233$. (b)~When the tickets are sampled with replacement, there are seven independent draws. Again we first find the probability of not winning a prize: $(29/30)^7 = 0.789$. Thus, the probability of winning (at least) one prize when drawing with replacement is 0.211.}

\begin{exercisewrap}
\begin{nexercise} \label{followUpToRaffleOf30TicketsWWOReplacement}
Compare your answers in Guided Practice~\ref{raffleOf30TicketsWWOReplacement}. How much influence does the sampling method have on your chances of winning a prize?\footnotemark
\end{nexercise}
\end{exercisewrap}
\footnotetext{There is about a 10\% larger chance of winning a prize when using sampling without replacement. However, at most one prize may be won under this sampling procedure.}

Had we repeated Guided Practice~\ref{raffleOf30TicketsWWOReplacement} with 300 tickets instead of 30, we would have found something interesting: the results would be nearly identical. The probability would be 0.0233 without replacement and 0.0231 with replacement. When the sample size is only a small fraction of the population (under 10\%), observations are nearly independent even when sampling without replacement.


{\input{ch_probability/TeX/sampling_from_a_small_population.tex}}





%_________________
\section{Random variables}
\label{randomVariablesSection}

\index{random variable|(}

\noindent%
It's often useful to model a process using what's
called a \term{random variable}.
Such a model allows us to apply a mathematical
framework and statistical principles for
better understanding and predicting outcomes
in the real world.

\begin{examplewrap}
\begin{nexample}{Two books are assigned for a statistics class: a textbook and its corresponding study guide. The university bookstore determined 20\% of enrolled students do not buy either book, 55\% buy the textbook only, and 25\% buy both books, and these percentages are relatively constant from one term to another. If~there are 100 students enrolled, how many books should the bookstore expect to sell to this class?}\label{bookStoreSales}
Around 20 students will not buy either book (0 books total), about 55 will buy one book (55 books total), and approximately 25 will buy two books (totaling 50 books for these 25 students). The bookstore should expect to sell about 105 books for this class.
\end{nexample}
\end{examplewrap}

\begin{exercisewrap}
\begin{nexercise}
Would you be surprised if the bookstore sold slightly more or less than 105 books?\footnotemark
\end{nexercise}
\end{exercisewrap}
\footnotetext{If they sell a little more or a little less, this should not be a surprise. Hopefully Chapter~\ref{introductionToData} helped make clear that there is natural variability in observed data. For example, if we would flip a coin 100 times, it will not usually come up heads exactly half the time, but it will probably be close.}

\begin{examplewrap}
\begin{nexample}{The textbook costs \$137 and the study guide \$33. How much revenue should the bookstore expect from this class of 100 students?}\label{bookStoreRev}
About 55 students will just buy a textbook, providing revenue of
\begin{align*}
\$137 \times  55 = \$7,535
\end{align*}
The roughly 25 students who buy both the textbook and the study guide would pay a total of
\begin{align*}
(\$137 + \$33) \times  25 = \$170 \times  25 = \$4,250
\end{align*}
Thus, the bookstore should expect to generate about $\$7,535 + \$4,250 = \$11,785$ from these 100 students for this one class. However, there might be some \emph{sampling variability} so the actual amount may differ by a little bit.
\end{nexample}
\end{examplewrap}

\begin{figure}[h]
  \centering
  \Figure[A probability distribution, which appears similar to a histogram. The horizontal axis is "Cost" and runs from \$0 to \$170. The vertical axis is Probability. There are three bars: a bar with height 0.2 at \$0, a bar with height 0.55 with height \$137, and a bar with height 0.25 at \$170. A red triangle is shown at the mean, located at \$117.85.]{0.6}{bookCostDist}
  \caption{Probability distribution for the bookstore's
      revenue from one student.
      The triangle represents
      the average revenue per student.}
  \label{bookCostDist}
\end{figure}

\D{\newpage}

\begin{examplewrap}
\begin{nexample}{What is the average revenue per student for this course?}\label{revFromStudent}
The expected total revenue is \$11,785, and there are 100 students. Therefore the expected revenue per student is $\$11,785/100 =  \$117.85$.
\end{nexample}
\end{examplewrap}

\subsection{Expectation}

\index{expectation|(}

We call a variable or process with a numerical outcome a \term{random variable}, and we usually represent this random variable with a capital letter such as $X$, $Y$, or $Z$. The amount of money a single student will spend on her statistics books is a random variable, and we represent it by $X$.

\begin{onebox}{Random variable}
A random process or variable with a numerical outcome.
\end{onebox}

The possible outcomes of $X$ are labeled with a corresponding lower case letter $x$ and subscripts. For example, we write $x_1=\$0$, $x_2=\$137$, and $x_3=\$170$, which occur with probabilities $0.20$, $0.55$, and $0.25$. The distribution of $X$ is summarized in Figure~\ref{bookCostDist} and Figure~\ref{statSpendDist}.

\begin{figure}[h]
\centering
\begin{tabular}{l ccc r}
\hline
$i$	  & 1 & 2 & 3  & Total\\
\hline
$x_i$ & \$0 & \$137 & \$170 & --\\
$P(X=x_i)$ & 0.20 & 0.55 & 0.25 & 1.00 \\
\hline
\end{tabular}
\caption{The probability distribution for the random variable $X$, representing the bookstore's revenue from a single student.}
\label{statSpendDist}
\end{figure}

We computed the average outcome of $X$ as \$117.85 in Example~\ref{revFromStudent}.
We call this average the \term{expected value} of $X$, denoted by $E(X)$\index{EX@$E(X)$}.
The expected value of a random variable is computed by adding each outcome weighted by its probability:
\begin{align*}
E(X) &= 0 \times  P(X=0) + 137 \times  P(X=137) + 170 \times  P(X=170) \\
	&= 0 \times  0.20 + 137 \times  0.55 + 170 \times  0.25 = 117.85
\end{align*}

\begin{onebox}{Expected value of a Discrete Random Variable}
If $X$ takes outcomes $x_1$, ..., $x_k$ with probabilities $P(X=x_1)$, ..., $P(X=x_k)$, the expected value of $X$ is the sum of each outcome multiplied by its corresponding probability:
\begin{align*}
E(X)
  &= x_1 \times P(X = x_1) + \cdots + x_k\times P(X = x_k) \\
  &= \sum_{i = 1}^{k} x_i P(X = x_i)
\end{align*}
The Greek letter $\mu$\index{Greek!mu@mu ($\mu$)}
may be used in place of the notation $E(X)$.
\end{onebox}

\D{\newpage}

The expected value for a random variable represents the average outcome. For example, $E(X)=117.85$ represents the average amount the bookstore expects to make from a single student, which we could also write as $\mu=117.85$.

It is also possible to compute the expected value of a continuous random variable (see Section~\ref{contDist}). However, it requires a little calculus and we save it for a later class.\footnote{$\mu = \int xf(x)dx$ where $f(x)$ represents a function for the density curve.}

In physics, the expectation holds the same meaning as the center of gravity. The distribution can be represented by a series of weights at each outcome, and the mean represents the balancing point. This is represented in Figures~\ref{bookCostDist} and~\ref{bookWts}. The idea of a center of gravity also expands to continuous probability distributions. Figure~\ref{contBalance} shows a continuous probability distribution balanced atop a wedge placed at the mean.

\begin{figure}
\centering
\Figure[A bar is hung by a string, and three weights are hanging on the bar at three different locations. The weights are located at the locations 0, 137, and 170 and have weights proportional to the probabilities 0.2, 0.55, and 0.25, respectively. The weights are balanced, because the string that is suspending the bar is located at the mean of the distribution, 117.85.]{0.72}{bookWts}
\caption{A weight system representing the probability distribution for $X$. The string holds the distribution at the mean to keep the system balanced.}
\label{bookWts}
\end{figure}

\begin{figure}
\centering
\Figure[A distribution that is skewed to the right is displayed, similar to how a histogram would appear if the bins were so small as to blend together and look continuous. This distribution is balancing atop a triangle located at the mean of the distribution.]{0.68}{contBalance}
\caption{A continuous distribution can also be balanced at its mean.}
\label{contBalance}
\end{figure}

\index{expectation|)}


\D{\newpage}

\subsection{Variability in random variables}

Suppose you ran the university bookstore. Besides how much revenue you expect to generate, you might also want to know the volatility (variability) in your revenue. 

The \indexthis{variance}{variance} and \indexthis{standard deviation}{standard deviation} can be used to describe the variability of a random variable. Section~\ref{variability}
introduced a method for finding the variance and standard deviation for a data set. We first computed deviations from the mean ($x_i - \mu$), squared those deviations, and took an average to get the variance. In the case of a random variable, we again compute squared deviations. However, we take their sum weighted by their corresponding probabilities, just like we did for the expectation. This weighted sum of squared deviations equals the variance, and we calculate the standard deviation by taking the square root of the variance, just as we did in Section~\ref{variability}.

\begin{onebox}{General variance formula}
If $X$ takes outcomes $x_1$, ..., $x_k$ with probabilities $P(X=x_1)$, ..., $P(X=x_k)$ and expected value $\mu=E(X)$, then the variance of $X$, denoted by $Var(X)$ or the symbol $\sigma^2$, is
\begin{align*}
\sigma^2 &= (x_1-\mu)^2\times P(X=x_1) + \cdots \\
	& \qquad\quad\cdots+ (x_k-\mu)^2\times P(X=x_k) \\
	&= \sum_{j=1}^{k} (x_j - \mu)^2 P(X=x_j)
\end{align*}
The standard deviation of $X$, labeled
$\sigma$\index{Greek!sigma@sigma ($\sigma$)},
is the square root of the variance.
\end{onebox}

\begin{examplewrap}
\begin{nexample}{Compute the expected value, variance, and standard deviation of $X$, the revenue of a single statistics student for the bookstore.}
It is useful to construct a table that holds computations for each outcome separately, then add up the results.
\begin{center}
\begin{tabular}{l rrr r}
\hline
$i$ & 1 & 2 & 3 & Total \\
\hline
$x_i$ & \$0 & \$137 & \$170 &  \\
$P(X=x_i)$ & 0.20 & 0.55 & 0.25 &  \\
$x_i \times  P(X=x_i)$ & 0 & 75.35 & 42.50 & 117.85 \\
\hline
\end{tabular}
\end{center}
Thus, the expected value is $\mu=117.85$, which we computed earlier. The variance can be constructed by extending this table:
\begin{center}
\begin{tabular}{l rrr r}
\hline
$i$ & 1 & 2 & 3 & Total \\
\hline
$x_i$ & \$0 & \$137 & \$170 &  \\
$P(X=x_i)$ & 0.20 & 0.55 & 0.25 &  \\
$x_i \times  P(X=x_i)$ & 0 & 75.35 & 42.50 & 117.85 \\
$x_i - \mu$ & -117.85 & 19.15 & 52.15 &  \\
$(x_i-\mu)^2$ & 13888.62 &  366.72 & 2719.62 &  \\
$(x_i-\mu)^2\times P(X=x_i)$ & 2777.7 & 201.7 & 679.9 & 3659.3 \\
\hline
\end{tabular}
\end{center}
The variance of $X$ is $\sigma^2 = 3659.3$, which means the standard deviation is $\sigma = \sqrt{3659.3} = \$60.49$.
\end{nexample}
\end{examplewrap}

\begin{exercisewrap}
\begin{nexercise}
The bookstore also offers a chemistry textbook for \$159 and a book supplement for \$41. From past experience, they know about 25\% of chemistry students just buy the textbook while 60\% buy both the textbook and supplement.\footnotemark
\begin{enumerate}
\item[(a)] What proportion of students don't buy either book? Assume no students buy the supplement without the textbook.
\item[(b)] Let $Y$ represent the revenue from a single student. Write out the probability distribution of $Y$, i.e. a table for each outcome and its associated probability.
\item[(c)] Compute the expected revenue from a single chemistry student. 
\item[(d)] Find the standard deviation to describe the variability associated with the revenue from a single student.
\end{enumerate}
\end{nexercise}
\end{exercisewrap}
\footnotetext{(a) 100\% - 25\% - 60\% = 15\% of students do not buy any books for the class. Part~(b) is represented by the first two lines in the table below. The expectation for part~(c) is given as the total on the line $y_i\times P(Y=y_i)$. The result of part~(d) is the square-root of the variance listed on in the total on the last line: $\sigma = \sqrt{Var(Y)} = \$69.28$.
\begin{center}
\begin{tabular}{rrrrr}
  \hline
$i$ (scenario) & 1 (\resp{noBook}) & 2 (\resp{textbook}) & 3 (\resp{both}) & Total \\
  \hline
$y_i$ & 0.00 & 159.00 & 200.00 &  \\
$P(Y=y_i)$ & 0.15 & 0.25 & 0.60 & \\
$y_i\times P(Y=y_i)$ & 0.00 & 39.75 & 120.00 & $E(Y) = 159.75$\\
$y_i-E(Y)$ & -159.75 & -0.75 & 40.25 & \\
$(y_i-E(Y))^2$ & 25520.06 & 0.56 & 1620.06 & \\
$(y_i-E(Y))^2\times P(Y)$ & 3828.0 & 0.1 & 972.0 & $Var(Y) \approx 4800$ \\
   \hline
\end{tabular}
\end{center}}

\subsection{Linear combinations of random variables}

So far, we have thought of each variable as being a complete story in and of itself. Sometimes it is more appropriate to use a combination of variables. For instance, the amount of time a person spends commuting to work each week can be broken down into several daily commutes. Similarly, the total gain or loss in a stock portfolio is the sum of the gains and losses in its components.

\begin{examplewrap}
\begin{nexample}{John travels to work five days a week. We will use $X_1$ to represent his travel time on Monday, $X_2$ to represent his travel time on Tuesday, and so on. Write an equation using $X_1$, ..., $X_5$ that represents his travel time for the week, denoted by $W$.}
His total weekly travel time is the sum of the five daily values:
\begin{align*}
W = X_1 + X_2 + X_3 + X_4 + X_5
\end{align*}
Breaking the weekly travel time $W$ into pieces provides a framework for understanding each source of randomness and is useful for modeling $W$.
\end{nexample}
\end{examplewrap}

\D{\newpage}

\begin{examplewrap}
\begin{nexample}{It takes John an average of 18 minutes each day to commute to work. What would you expect his average commute time to be for the week?}
We were told that the average (i.e. expected value) of the commute time is 18 minutes per day: $E(X_i) = 18$. To get the expected time for the sum of the five days, we can add up the expected time for each individual day:
\begin{align*}
E(W) &= E(X_1 + X_2 + X_3 + X_4 + X_5) \\
	&= E(X_1) + E(X_2) + E(X_3) + E(X_4) + E(X_5) \\
	&= 18 + 18 + 18 + 18 + 18 = 90\text{ minutes}
\end{align*}
The expectation of the total time is equal to the sum of the expected individual times. More generally, the expectation of a sum of random variables is always the sum of the expectation for each random variable.
\end{nexample}
\end{examplewrap}

\begin{exercisewrap}
\begin{nexercise}
\label{elenaIsSellingATVAndBuyingAToasterOvenAtAnAuction}%
Elena is selling a TV at a cash auction and also intends to buy a toaster oven in the auction. If $X$ represents the profit for selling the TV and $Y$ represents the cost of the toaster oven, write an equation that represents the net change in Elena's cash.\footnotemark
\end{nexercise}
\end{exercisewrap}
\footnotetext{She will make $X$ dollars on the TV but spend $Y$ dollars on the toaster oven: $X-Y$.}

\begin{exercisewrap}
\begin{nexercise}
Based on past auctions, Elena figures she should expect to make about \$175 on the TV and pay about \$23 for the toaster oven. In total, how much should she expect to make or spend?\footnotemark
\end{nexercise}
\end{exercisewrap}
\footnotetext{$E(X-Y) = E(X) - E(Y) = 175 - 23 = \$152$. She should expect to make about \$152.}

\begin{exercisewrap}
\begin{nexercise} \label{explainWhyThereIsUncertaintyInTheSum}
Would you be surprised if John's weekly commute wasn't exactly 90 minutes or if Elena didn't make exactly \$152? Explain.\footnotemark
\end{nexercise}
\end{exercisewrap}
\footnotetext{No, since there is probably some variability. For example, the traffic will vary from one day to next, and auction prices will vary depending on the quality of the merchandise and the interest of the attendees.}

Two important concepts concerning combinations of random variables have so far been introduced. First, a final value can sometimes be described as the sum of its parts in an equation. Second, intuition suggests that putting the individual average values into this equation gives the average value we would expect in total. This second point needs clarification -- it is guaranteed to be true in what are called \emph{linear combinations of random variables}.

A \term{linear combination} of two random variables $X$ and $Y$ is a fancy phrase to describe a combination
\begin{align*}
aX + bY
\end{align*}
where $a$ and $b$ are some fixed and known numbers. For John's commute time, there were five random variables -- one for each work day -- and each random variable could be written as having a fixed coefficient of 1:
\begin{align*}
1X_1 + 1 X_2 + 1 X_3 + 1 X_4 + 1 X_5
\end{align*}
For Elena's net gain or loss, the $X$ random variable had
a coefficient of +1 and the $Y$ random variable had
a coefficient of~-1.

\D{\newpage}

When considering the average of a linear combination of random variables, it is safe to plug in the mean of each random variable and then compute the final result. For a few examples of nonlinear combinations of random variables -- cases where we cannot simply plug in the means -- see the footnote.\footnote{If $X$ and $Y$ are random variables, consider the following combinations: $X^{1+Y}$, $X\times Y$, $X/Y$. In such cases, plugging in the average value for each random variable and computing the result will not generally lead to an accurate average value for the end result.}

\begin{onebox}{Linear combinations of random variables and the average result}
If $X$ and $Y$ are random variables, then a linear combination of the random variables is given by
\begin{align*}
aX + bY
\end{align*}
where $a$ and $b$ are some fixed numbers. To compute the average value of a linear combination of random variables, plug in the average of each individual random variable and compute the result:
\begin{align*}
a\times E(X) + b\times E(Y)
\end{align*}
Recall that the expected value is the same as the mean, e.g. $E(X) = \mu_X$.
\end{onebox}

\begin{examplewrap}
\begin{nexample}{Leonard has invested \$6000 in Caterpillar Inc
    (stock ticker: CAT) and \$2000 in Exxon Mobil Corp (XOM).
    If $X$ represents the change in Caterpillar's stock next month
    and $Y$ represents the change in Exxon Mobil's stock
    next month, write an equation that describes how much
    money will be made or lost in Leonard's stocks for the
    month.}
  For simplicity, we will suppose $X$ and $Y$ are not
  in percents but are in decimal form
  (e.g. if Caterpillar's stock increases 1\%, then $X=0.01$;
  or if it loses 1\%, then $X=-0.01$).
  Then we can write an equation for Leonard's gain as
  \begin{align*}
  \$6000\times X + \$2000\times Y
  \end{align*}
  If we plug in the change in the stock value for $X$ and $Y$,
  this equation gives the change in value of Leonard's stock
  portfolio for the month. A positive value represents a gain,
  and a negative value represents a loss.
\end{nexample}
\end{examplewrap}

\begin{exercisewrap}
\begin{nexercise}\label{expectedChangeInLeonardsStockPortfolio}
Caterpillar stock has recently been rising
at 2.0\% and Exxon Mobil's at 0.2\% per month, respectively.
Compute the expected change in Leonard's stock portfolio
for next month.\footnotemark
\end{nexercise}
\end{exercisewrap}
% library(openintro); d <- stocks_18; cols <- 2:4; apply(d[, cols], 2, mean); apply(d[, cols], 2, sd)
\footnotetext{%
  $E(\$6000\times X + \$2000\times Y) =
    \$6000\times 0.020 + \$2000\times 0.002 = \$124$.}

\begin{exercisewrap}
\begin{nexercise}
You should have found that Leonard expects a positive gain
in Guided Practice~\ref{expectedChangeInLeonardsStockPortfolio}.
However, would you be surprised if he actually had
a loss this month?\footnotemark
\end{nexercise}
\end{exercisewrap}
\footnotetext{No.
  While stocks tend to rise over time, they are often
  volatile in the short term.}


\D{\newpage}

\subsection{Variability in linear combinations of random variables}
\label{var_lin_combo_of_RVs}

Quantifying the average outcome from a linear combination
of random variables is helpful, but it is also important
to have some sense of the uncertainty associated with
the total outcome of that combination of random variables.
The expected net gain or loss of Leonard's stock portfolio
was considered in Guided Practice~\ref{expectedChangeInLeonardsStockPortfolio}.
However, there was no quantitative discussion of the
volatility of this portfolio.
For instance, while the average monthly gain might be
about \$124 according to the data, that gain is not guaranteed.
Figure~\ref{changeInLeonardsStockPortfolioFor36Months}
shows the monthly changes in a portfolio like Leonard's during
a three year period.
The gains and losses vary widely, and quantifying these
fluctuations is important when investing in stocks.

\begin{figure}[ht]
\centering
\Figure[A dot plot is overlaid on a box plot for a variable "Monthly Returns Over 3 Years". The box portion spans about -200 to 450, with the median line at about 200. The whiskers extend to the lower end at about -600 and at the upper end at about 1050. There is one point beyond the lower whisker located at about -1400. The the dot plot, the points are pretty evenly distributed across the locations of the box-and-whiskers portion of the box plot, with the one exception being the point at -1400.]{0.6}{changeInLeonardsStockPortfolioFor36Months}
\caption{The change in a portfolio like Leonard's for 36 months,
    where \$6000 is in Caterpillar's stock and \$2000 is in
    Exxon Mobil's.}
\label{changeInLeonardsStockPortfolioFor36Months}
\end{figure}

Just as we have done in many previous cases,
we use the variance and standard deviation to describe
the uncertainty associated with Leonard's monthly returns.
To do so, the variances of each stock's monthly return
will be useful, and these are shown in
Figure~\ref{sumStatOfCATXOM}.
The stocks' returns are nearly independent.

\begin{figure}
\centering
\begin{tabular}{lrrr}
\hline
    & Mean ($\bar{x}$) & Standard deviation ($s$) &
	    Variance ($s^2$) \\
\hline
CAT & 0.0204  & 0.0757  & 0.0057 \\
XOM & 0.0025  & 0.0455  & 0.0021 \\
\hline
\end{tabular}
\caption{The mean, standard deviation, and variance of the
    CAT and XOM stocks.
    These statistics were estimated from historical
    stock data, so notation used for sample statistics
    has been used.}
\label{sumStatOfCATXOM}
\end{figure}

Here we use an equation from probability theory to
describe the uncertainty of Leonard's monthly returns;
we leave the proof of this method to a dedicated
probability course.
The variance of a linear combination of random variables
can be computed by plugging in the variances of the
individual random variables and squaring the coefficients
of the random variables:
\begin{align*}
Var(aX + bY) = a^2\times Var(X) + b^2\times Var(Y)
\end{align*}
It is important to note that this equality assumes the
random variables are independent;
%\Comment{new description here about if independence is broken}
if independence doesn't hold, then a modification to
this equation would be required that we leave as a topic
for a future course to cover.
This equation can be used to compute the variance of
Leonard's monthly return:
\begin{align*}
Var(6000\times X + 2000\times Y)
	&= 6000^2\times Var(X) + 2000^2\times Var(Y) \\
	&= 36,000,000\times 0.0057 + 4,000,000\times 0.0021 \\
	&\approx 213,600
% sum(c(36e6, 4e6) * c(0.0057, 0.0021))
\end{align*}
The standard deviation is computed as the square root
of the variance: $\sqrt{213,600} = \$463$.
While an average monthly return of \$124 on an
\$8000 investment is nothing to scoff at,
the monthly returns are so volatile that Leonard should
not expect this income to be very stable.

\begin{onebox}{Variability of linear combinations of random variables}
The variance of a linear combination of random variables may be computed by squaring the constants, substituting in the variances for the random variables, and computing the result:
\begin{align*}
Var(aX + bY) = a^2\times Var(X) + b^2\times Var(Y)
\end{align*}
This equation is valid as long as the random variables are independent of each other. The standard deviation of the linear combination may be found by taking the square root of the variance.
\end{onebox}

\begin{examplewrap}
\begin{nexample}{Suppose John's daily commute has a standard deviation of 4 minutes. What is the uncertainty in his total commute time for the week?} \label{sdOfJohnsCommuteWeeklyTime}
The expression for John's commute time was
\begin{align*}
X_1 + X_2 + X_3 + X_4 + X_5
\end{align*}
Each coefficient is 1, and the variance of each day's time is $4^2=16$. Thus, the variance of the total weekly commute time is
\begin{align*}
&\text{variance }= 1^2 \times  16 + 1^2 \times  16 + 1^2 \times  16 + 1^2 \times  16 + 1^2 \times  16 = 5\times 16 = 80 \\
&\text{standard deviation } = \sqrt{\text{variance}} = \sqrt{80} = 8.94
\end{align*}
The standard deviation for John's weekly work commute time is about 9 minutes.
\end{nexample}
\end{examplewrap}

\begin{exercisewrap}
\begin{nexercise}
The computation in Example~\ref{sdOfJohnsCommuteWeeklyTime} relied on an important assumption: the commute time for each day is independent of the time on other days of that week. Do you think this is valid? Explain.\footnotemark
\end{nexercise}
\end{exercisewrap}
\footnotetext{One concern is whether traffic patterns tend to have a weekly cycle (e.g. Fridays may be worse than other days). If that is the case, and John drives, then the assumption is probably not reasonable. However, if John walks to work, then his commute is probably not affected by any weekly traffic cycle.}

\begin{exercisewrap}
\begin{nexercise}\label{elenaIsSellingATVAndBuyingAToasterOvenAtAnAuctionVariability}
Consider Elena's two auctions from Guided Practice~\ref{elenaIsSellingATVAndBuyingAToasterOvenAtAnAuction} on page~\pageref{elenaIsSellingATVAndBuyingAToasterOvenAtAnAuction}. Suppose these auctions are approximately independent and the variability in auction prices associated with the TV and toaster oven can be described using standard deviations of \$25 and \$8. Compute the standard deviation of Elena's net gain.\footnotemark
\end{nexercise}
\end{exercisewrap}
\footnotetext{The equation for Elena can be written as
\begin{align*}
(1)\times X + (-1)\times Y
\end{align*}
The variances of $X$ and $Y$ are 625 and 64. We square the coefficients and plug in the variances:
\begin{align*}
(1)^2\times Var(X) + (-1)^2\times Var(Y) = 1\times 625 + 1\times 64 = 689
\end{align*}
The variance of the linear combination is 689, and the standard deviation is the square root of 689: about \$26.25.}

Consider again Guided Practice~\ref{elenaIsSellingATVAndBuyingAToasterOvenAtAnAuctionVariability}. The negative coefficient for $Y$ in the linear combination was eliminated when we squared the coefficients. This generally holds true: negatives in a linear combination will have no impact on the variability computed for a linear combination, but they do impact the expected value computations.

\index{random variable|)}


{\input{ch_probability/TeX/random_variables.tex}}





%_________________
\section{Continuous distributions}
\label{contDist}

\noindent%
So far in this chapter we've discussed cases
where the outcome of a variable is discrete.
In this section, we consider a context where
the outcome is a continuous numerical variable.

\index{data!US adult heights|(}
\index{hollow histogram|(}
\begin{examplewrap}
\begin{nexample}{Figure~\ref{fdicHistograms} shows a few
    different hollow histograms for the heights of US adults.
    How does changing the number of bins allow you to make
    different interpretations of the data?}
  \label{usHeights}%
  Adding more bins provides greater detail.
  This sample is extremely large, which is why much smaller
  bins still work well.
  Usually we do not use so many bins with smaller sample
  sizes since small counts per bin mean the bin heights
  are very volatile.
\end{nexample}
\end{examplewrap}

\begin{figure}[ht]
  \centering
  \Figure[Four hollow histograms are shown for the US adult heights in centimeters with varying bin widths. The difference in appears will first be discussed, and then the shape of the last, most detailed histogram will be given. The first histogram has about 6 bins with values that appear to be non-zero, so the outline is very boxy. The second histogram has about 12 non-zero bins, and appears a bit more continuous and less boxy than the first histogram. The third histogram has about 25 non-zero bins, and the hollow histogram outline is starting to look somewhat smoother. The last histogram has about 50 bins, and due to the large number of bins, the distribution looks quite smooth, in that no steps from one bin to the next is a substantial jump or drop in height. Next, this last histogram is described: The bin heights are about zero until 147, then they steadily climb up to about 155 before steeply climbing a little until 157 and then steadily climb to a peak at about 165. From here the histogram declines about 10\% from its peak at 170, at which point the decline is more gradual until about 183, at which point it descends rapidly until about 187 where it begins to descend more slowly as it approaches 0. At about 200, the bin heights have essentially hit zero and stay there.]{}{fdicHistograms}
  \caption{Four hollow histograms of US adults heights
      with varying bin widths.}
  \label{fdicHistograms}
\end{figure}

\begin{examplewrap}
\begin{nexample}{What proportion of the sample is between \resp{180} cm and \resp{185} cm tall (about 5'11" to 6'1")?}\label{contDistProb}
We can add up the heights of the bins in the range \resp{180} cm and \resp{185} and divide by the sample size. For instance, this can be done with the two shaded bins shown in Figure~\ref{usHeightsHist180185}. The two bins in this region have counts of 195,307 and 156,239 people, resulting in the following estimate of the probability:
\begin{align*}
\frac{195307 + 156239}{\text{3,000,000}} = 0.1172
\end{align*}
This fraction is the same as the proportion of the histogram's area that falls in the range \resp{180} to \resp{185} cm.
\end{nexample}
\end{examplewrap}

\begin{figure}[h]
  \centering
  \Figure[A histogram for heights is shown, with the two histogram bins between 180 and 185 centimeters are shaded, representing the individuals with heights between 180 and 185 centimeters.]{0.9}{usHeightsHist180185}
  \caption{A histogram with bin sizes of 2.5 cm.
      The shaded region represents individuals with
      heights between \resp{180} and \resp{185} cm.}
  \label{usHeightsHist180185}
\end{figure}


\D{\newpage}

\subsection{From histograms to continuous distributions}

Examine the transition from a boxy hollow histogram in the top-left of Figure~\ref{fdicHistograms} to the much smoother plot in the lower-right. In this last plot, the bins are so slim that the hollow histogram is starting to resemble a smooth curve. This suggests the population height as a \emph{continuous} numerical variable might best be explained by a curve that represents the outline of extremely slim bins.

This smooth curve represents a
\termsub{probability density function}
    {probability!density function}
(also called a \term{density} or \term{distribution}), and such a curve is shown in Figure~\ref{fdicHeightContDist} overlaid on a histogram of the sample. A density has a special property: the total area under the density's curve is 1. 

\begin{figure}[tbh]
\centering
\Figure[A histogram for heights of US adults is shown with an overlaid continuous line along the heights of the bins. This continuous line is smooth and would represent what a hollow histogram would look like if we had infinite data and the bin width was so small that the boxy outline of the histogram looks continuous. This is called a "continuous probability density".]{0.9}{fdicHeightContDist}
\caption{The continuous probability distribution of heights for US adults.}
\label{fdicHeightContDist}
\end{figure}

\index{hollow histogram|)}


\D{\newpage}

\subsection{Probabilities from continuous distributions}

We computed the proportion of individuals with heights \resp{180} to \resp{185} cm in Example~\ref{contDistProb} as a fraction:
\begin{align*}
\frac{\text{number of people between \resp{180} and \resp{185}}}{\text{total sample size}}
\end{align*}
We found the number of people with heights between \resp{180} and \resp{185} cm by determining the fraction of the histogram's area in this region. Similarly, we can use the area in the shaded region under the curve to find a probability (with the help of a computer):
\begin{align*}
P(\text{\var{height} between \resp{180} and \resp{185}})
	= \text{area between \resp{180} and \resp{185}}
	= 0.1157
\end{align*}
The probability that a randomly selected person is between \resp{180} and \resp{185} cm is 0.1157. This is very close to the estimate from Example~\ref{contDistProb}: 0.1172. 

\begin{figure}[h]
  \centering
  \Figure[A density curve for heights is shown, with the region between 180 and 185 centimeters being shaded.]{0.7}{fdicHeightContDistFilled}
  \caption{Density for heights in the US adult population
      with the area between 180 and 185 cm shaded.
      Compare this plot with Figure~\ref{usHeightsHist180185}.}
  \label{fdicHeightContDistFilled}
\end{figure}

\begin{exercisewrap}
\begin{nexercise}
Three US adults are randomly selected. The probability a single adult is between \resp{180} and \resp{185} cm is 0.1157.\footnotemark\vspace{-1.5mm}
\begin{enumerate}
\setlength{\itemsep}{0mm}
\item[(a)] What is the probability that all three are between \resp{180} and \resp{185} cm tall?
\item[(b)] What is the probability that none are between \resp{180} and \resp{185} cm?
\end{enumerate}
\end{nexercise}
\end{exercisewrap}
\footnotetext{Brief answers:
  (a) $0.1157 \times 0.1157 \times 0.1157 = 0.0015$.
  (b) $(1-0.1157)^3 = 0.692$}

\begin{examplewrap}
\begin{nexample}{What is the probability that a randomly selected person is \textbf{exactly} \resp{180}~cm? Assume you can measure perfectly.}
\label{probabilityOfExactly180cm}
This probability is zero. A person might be close to \resp{180} cm, but not exactly \resp{180} cm tall. This also makes sense with the definition of probability as area; there is no area captured between \resp{180}~cm and \resp{180}~cm.
\end{nexample}
\end{examplewrap}

\begin{exercisewrap}
\begin{nexercise}
Suppose a person's height is rounded to the nearest centimeter. Is there a chance that a random person's \textbf{measured} height will be \resp{180} cm?\footnotemark
\end{nexercise}
\end{exercisewrap}
\footnotetext{This has positive probability. Anyone between \resp{179.5} cm and \resp{180.5} cm will have a \emph{measured} height of \resp{180} cm. This is probably a more realistic scenario to encounter in practice versus Example~\ref{probabilityOfExactly180cm}.}

\index{data!US adult heights|)}


{\input{ch_probability/TeX/continuous_distributions.tex}}
